%%%%%%%%%%%%%%%%%%%%%%%%%%%%%%%%%%%%%%%%%
%%            LMU-Vorlage              %%
%%                                     %%
%%         zur Erstellung einer        %%
%%   Dissertation mit pdflatex/latex   %%
%%                                     %%
%%  (2002) Robert Dahlke               %%
%%         & Sigmund Stintzing         %%
%%%%%%%%%%%%%%%%%%%%%%%%%%%%%%%%%%%%%%%%%

\documentclass[12pt]{book}


%%%%%%%%%%%%%%%%%%%%%%%%%%%%
%%   Zusaetzliche Pakete  %%
%%%%%%%%%%%%%%%%%%%%%%%%%%%%

\usepackage{a4wide}
\usepackage[utf8]{inputenc}
\usepackage[T1]{fontenc}
\usepackage{fancyhdr}
\usepackage{graphicx}
\usepackage{caption}
\usepackage[toc,page]{appendix}
\usepackage[bookmarks=true,draft=false]{hyperref}
\usepackage[dvipsnames]{xcolor}
\usepackage{wrapfig}
\usepackage{amsmath}
\usepackage{amssymb}
\usepackage{amsthm}
\usepackage{dsfont}
\usepackage{mathtools}
\usepackage{relsize}
\usepackage{algorithm}
\usepackage{algpseudocode}
\usepackage{colonequals}
\usepackage{units}
\usepackage{mymacros}

\usepackage{listings}

\usepackage[nottoc]{tocbibind}

\usepackage[backend=bibtex]{biblatex}
\addbibresource{references.bib}

\usepackage{mymacros}

\usepackage[toc,xindy,nonumberlist]{glossaries}
\makeglossaries
\loadglsentries{glossary}


%%%%%%%%%%%%%%%%%%%%%%%%%%%
%% Definition of Heading %%
%%%%%%%%%%%%%%%%%%%%%%%%%%%

\pagestyle{fancyplain}
\renewcommand{\chaptermark}[1]%
         {\markboth{\thechapter.\ #1}{}}
\renewcommand{\sectionmark}[1]%
         {\markright{\thesection\ #1}}
\lhead[\fancyplain{}{\bfseries\thepage}]%
    {\fancyplain{}{\bfseries\rightmark}}
\rhead[\fancyplain{}{\bfseries\leftmark}]%
    {\fancyplain{}{\bfseries\thepage}}
\cfoot{}


%%%%%%%%%%%%%%%%%%%%%%%%%%%%%%%%%%%%%%%%%%%%%%%%%%%%%
%%  Definition des Deckblattes und der Titelseite  %%
%%%%%%%%%%%%%%%%%%%%%%%%%%%%%%%%%%%%%%%%%%%%%%%%%%%%%

\newcommand{\LMUTitle}[9]{
  \thispagestyle{empty}
  \vspace*{\stretch{1}}
  {\parindent0cm
   \rule{\linewidth}{.7ex}}
  \begin{flushright}

    \vspace*{\stretch{1}}
    \sffamily\bfseries\Huge
    #1\\
    \vspace*{\stretch{1}}
    \sffamily\bfseries\large
    #2
    \vspace*{\stretch{1}}
  \end{flushright}
  \rule{\linewidth}{.7ex}
  \vspace*{\stretch{5}}
  \begin{center}
    \includegraphics[width=2in]{siegel}
  \end{center}
  \vspace*{\stretch{1}}
  \begin{center}\sffamily\LARGE{#5}\end{center}
  \newpage
  \thispagestyle{empty}

  \cleardoublepage
  \thispagestyle{empty}

  \vspace*{\stretch{1}}
  {\parindent0cm
  \rule{\linewidth}{.7ex}}
  \begin{flushright}
    \vspace*{\stretch{1}}
    \sffamily\bfseries\Huge
    #1\\
    \vspace*{\stretch{1}}
    \sffamily\bfseries\large
    #2
    \vspace*{\stretch{1}}
  \end{flushright}
  \rule{\linewidth}{.7ex}

  \vspace*{\stretch{3}}
  \begin{center}
    \Large Masterarbeit\\
    \Large an der #4\\
    \Large der Ludwig--Maximilians--Universität\\
    \Large M\"unchen\\
    \vspace*{\stretch{1}}
    \Large vorgelegt von\\
    \Large #2\\
    \Large aus #3\\
    \vspace*{\stretch{2}}
    \Large M\"unchen, den #6
  \end{center}

  \newpage
  \thispagestyle{empty}

  \vspace*{\stretch{1}}

  \begin{flushleft}
    \large Erstgutachter:  #7 \\[1mm]
    \large Diese Masterarbeit steht unter einer \ccby-Lizenz \\
  \end{flushleft}

  \cleardoublepage
}

%%%%%%%%%%%%%%%%%%%%%%%%%%%%
%%  Beginn des Dokuments  %%
%%%%%%%%%%%%%%%%%%%%%%%%%%%%

\begin{document}


  \frontmatter


  \LMUTitle
      {The de Moivre-Laplace theorem}                  % Titel der Arbeit
      {Stephan Kulla}                                  % Vor- und Nachname des Autors
      {Potsdam}                                        % Geburtsort des Autors
      {Mathematikfakultät}                             % Name der Fakultaet
      {München 2015}                                   % Ort und Jahr der Erstellung
      {15.02.2013}                                     % Tag der Abgabe
      {Prof. Dr. Peter Pickl}                          % Name des Erstgutachters

  %\markboth{Summary}{Summary}
  %\include{summary}
  %\cleardoublepage

  \tableofcontents
  \markboth{Contents}{Contents} 
  \cleardoublepage

  \glsaddall
  \printglossary[title={List of symbols},style=altlist]
  \cleardoublepage

 
  %\markboth{Remarks for this thesis}{Remarks for this thesis}
  %\include{remarks}

  \mainmatter\setcounter{page}{1}
  \chapter{Introduction}

In this thesis, I present a new proof for the theorem by de Moivre and Laplace. This theorem states that the standardized binomial distribution given by the probability mass function
\begin{align}
  \bb{\x} = \binom nk p^k q^{n-k}
\end{align}
with $\x[k] = \frac{k-np}{\sqrt{npq}}$ for $p\in (0,1)$, $q=1-p$, $k,n\in\N$, and $k\le n$ can be approximated by the standardized normal distribution, whose density function is $\fphi{x} = \frac{1}{\sqrt {2\pi}} \fexp{-\frac{x^2}2}$.

The new proof approximates the slope between neighboring $\bb{\x}$ and derives from it an ordinary differential equation (ODE) whose solution is the normal distribution. Hence it will give an interesting application of the theory of ODEs (especially one-step methods) to probability theory. There are two possible advantages of the new proof which we will investigate in this thesis:

\begin{enumerate}
  \item I will compare the error estimate of the standard proof using Stirling's formula, with the error bounds provided by the new proof. Thus, we can see whether the new proof yield better approximations or better error estimations.

  \item The new proof may be easier to understand than the standard proof, and thus may be used in a course or a textbook.
\end{enumerate}

  \chapter{Asymptotic notations} \label{chapter:tools}

In this chapter I will introduce and explain the asymptotic notations like $\bigo{\cdot}$ and $\littleo{\cdot}$ we will use in this thesis.

\section{The big O notation}

To describe the asymptotic behavior of an error, the big O notation can be utilized. It declares the convergence speed with which the error converges at least to zero and means~\cite[p. 444]{graham}\cite[p.~100]{aigner}

\begin{definition}[Big O notation]
  Let $\seq{a_n}$ be a sequence with positive members. $\bigo{a_n}$ is the set of all sequences defined via
  \begin{align} \label{bigo_def1}
    \seq{\epsilon_n} \in \bigo{a_n} :\iff \exists C > 0\, \forall n\in \N: \abs{\epsilon_n} \le C \cdot \abs{a_n}
  \end{align}
\end{definition}

\noindent This definition is equivalent to~\cite[p.~383]{hachenberger}\cite{wiki:bigo}
\begin{align} \label{bigo_def_limsup}
  \seq{\epsilon_n} \in \bigo{a_n} :\iff \exists C_\infty > 0 \colon\limsup_{n\to\infty} \abs{\frac{\epsilon_n}{a_n}} \le C_\infty
\end{align}

So $\epsilon_n \in \bigo{\frac 1n}$ means that $\epsilon_n$ approaches zero at least with the convergence speed of $\seq{\frac 1n}$. Thereby the advantage of the big O notation is its simple arithmetic rules. For example we have
\begin{align}
  \bigo{\frac 1n} + \bigo{\frac 1{n^2}} \subseteq \bigo{\frac 1n}
\end{align}
These arithmetic rules (which we will discuss in the end of this section) simplify error estimations.

\section{The big Delta and the big Psi notation}
Unfortunately the big O notation does only provide an estimate of the convergence speed but no estimation of the error's value since the constants $C$ and $C_\infty$ in \eqref{bigo_def1} and \eqref{bigo_def_limsup} are not known\footnote{cf.~\cite[p.~444]{graham},~\cite{hurkyl_bigo} and \cite{templatetypedef_bigo}}. Therefore I want to introduce the big Delta and the big Psi notation with which besides to the convergence speed also an estimate for the absolute error and the asymptotic error can be stated. Both notations are defined as\footnote{Donald Knuth uses the term ``A notation'' for the big Delta notation \cite{anotation}. Please compare also \cite{tampis_bigpsi} and~\cite{tampis_bigabs} respectively for the usage of the big Psi notation and big Delta notation in mathematics.}:

\begin{definition}[Big Delta notation for the absolute error]
  Let $\seq{a_n}$ be a sequence. $\bigabs{a_n}$ is the set of all sequences $\seq{\epsilon_n}$ with $\abs{\epsilon_n} \le \abs{a_n}$ for all $n\in\N$:
  \begin{align}
    \seq{\epsilon_n} \in \bigabs{a_n} \iff \forall n\in\N: \abs{\epsilon_n} \le \abs{a_n}
  \end{align}
\end{definition}

\begin{definition}[Big Psi notation for the asymptotic error]
  Let $\seq{a_n}$ be a sequence with $a_n > 0$ for all $n\in\N$. $\bigpsi{a_n}$ is the set of all sequences $\seq{\epsilon_n}$ fulfilling $\limsup_{n\to\infty} \abs{\frac{\epsilon_n}{a_n}} \le 1$. Thus
  \begin{align}
    \seq{\epsilon_n} \in \bigpsi{a_n} :\iff \limsup_{n\to\infty} \abs{\frac{\epsilon_n}{a_n}} \le 1
  \end{align}
\end{definition}

Note that the definition of the big Psi notation is the same as the limes superior definition \eqref{bigo_def_limsup} of the big O notation whereby $C_\infty$ has the concrete value $1$ and the big Delta notation is the same as \eqref{bigo_def1} with $C=1$. Thus we directly see that both notations are specializations of the big O notation:
\begin{align}
  \seq{\epsilon_n} \in \bigabs{a_n} & \implies \seq{\epsilon_n} \in \bigo{a_n} \nl
  \seq{\epsilon_n} \in \bigpsi{a_n} & \implies \seq{\epsilon_n} \in \bigo{a_n}
\end{align}

\section{The little o notation}

We will also use the little o notation which is defined as~\cite[pp.~99,~103]{aigner}\cite[p.~385]{hachenberger}\cite{wiki:bigo}

\begin{definition}[Little o notation]
  Let $\seq{a_n}$ be a sequence with positive members\footnote{To allow $a_n=0$ it is possible to define the little o notation via~\cite[pp.~448]{graham}\cite{wiki:bigo}
  \begin{align}
    \seq{b_n}\in\littleo{a_n} :\iff \forall \epsilon > 0\, \exists N\in\N\, \forall n \ge N: \abs{b_n} \le \epsilon \abs{a_n}
  \end{align}
   In this work we will only deal with sequences $\seq{a_n}$ with $a_n > 0$ for all $n\in\N$.
}. $\littleo{a_n}$ is the set of all sequences $\seq{\epsilon_n}$ with $\lim_{n\to\infty} \abs{\frac{\epsilon_n}{a_n}} = 0$. Thus
  \begin{align}
    \seq{\epsilon_n}\in\littleo{a_n} \iff \lim_{n\to\infty} \abs{\frac{\epsilon_n}{a_n}} = 0
  \end{align}
\end{definition}

\section{Calculations with sets} \label{sec:calculations_sets}

With the above notations we can simplify error estimations. Instead of stating a concrete error $\epsilon_n$ of an approximation at the $n$-th step, we use one of the above notations in the form $S(a_n)$ with $S\in\{O, \Delta, \Psi, o\}$. By using $S(a_n)$ in an expression we mean that there is a $\seq{\epsilon_n}\in S(a_n)$ fulfilling the statement. Because we will have to deal with expressions containing sets of sequences, we need rules for interpreting such expressions. These rules are \cite{bigo_rules1}\cite[pp. 17-18]{kulla}:

\begin{enumerate}
  \item Let $f:D\subseteq\R^n\to\R$ be an $n$-ary operation of real numbers and let $A_1,A_2,\ldots,A_n$ be $n$ subsets of $\R$ such that $A_1\times A_2\times \dots\times A_n\subseteq D$. The result  $f\left(A_1,\ldots,A_n\right)$ of the operation with all $n$ sets is defined as the image of $f$ \footnote{The extended version of $f$ on sets calls Ramon E. Moore in his book about interval analysis the ``unit extension''~\cite[p. 18]{moore}}
\begin{align}
  f\left(A_1,\ldots A_n\right)=\left\{f(x_1,\ldots,x_n):x_1\in A_1,x_2\in A_2,\ldots,x_n\in A_n\right\}
\end{align}
  For example
  \begin{align}
    \bigo{a_n} + \bigabs{b_n} = \{ \seq{\epsilon_n+\delta_n} : \seq{\epsilon_n}\in \bigo{a_n} \land \seq{\delta_n}\in\bigabs{b_n}\}
  \end{align}

  \item Now consider a function $\tilde f:D\subseteq\R^n\to \mathcal P(\R)$ which maps $n$ numbers to a set of real numbers. The result of $\tilde f\left(A_1,A_2,\ldots,A_n\right)$ for $A_1\times\dots\times A_n\subseteq D$ shall be
\begin{align}
  \tilde f\left(A_1,A_2,\ldots,A_n\right) = \bigcup_{x_1\in A_1, \ldots, x_n \in A_n} f(x_1,\ldots,x_n)
\end{align}
Thus $\tilde f(A_1,\ldots,A_n)$ is the union of all possible images of $\tilde f$. For example
\begin{align}
  \littleo{\bigo{a_n}} = \bigcup_{\seq{\epsilon_n}\in\bigo{a_n}} \littleo{\epsilon_n}
\end{align}

  \item If in the above rules one of the $A_k$ is a sequence $\epsilon_n$ (and not a set), we have to interpret $A_k$ as the singleton $\{\seq{\epsilon_n}\}$. Exceptions are expressions such as $\bigo{\epsilon_n}$ which are defined as in the above sections.
\end{enumerate}

It is also possible to use the rules stated by Henning Makholm\cite{bigo_rules2} (which he only formulated for the big O notation):

\begin{quote}
  ``Whenever you see an equation $t=u$ where $t$ or $u$ contain one or more $\bigo{\ldots}$, you're supposed to mentally do the following expansion:
  \begin{itemize}
    \item Replace \emph{each textual occurrence} of $O(\cdots)$ with $\phi(n)$ where $\phi$ is a \emph{fresh} function letter that ranges over $\mathcal O(\cdots)$. Different [occurrences] of $O(\cdots)$ in the formula should get separate function letters, even if the "$\cdots$" are the same.
    \item Every fresh $\phi$ that arises from an $O(\cdots)$ on the \emph{left} side of the equals sign should be \emph{universally} quantified.
    \item Every fresh $\phi$ that arises from an $O(\cdots)$ on the \emph{right} side of the equals sign should be \emph{existentially} quantified.
    \item The universal quantifiers on $\phi$s come before the existential ones.
    \item The variable $n$ that tends to a limit should be universally quantified \emph{after} all of the quantifiers on $\phi$s.''
  \end{itemize}
\end{quote}

\begin{remark}
  Be aware that I will use variables like $h=\frac{1}{\sqrt{npq}}$ where I drop the dependence on $n$ in the name of the variable. When I state something like $\littleo{h}$ I will always mean $\littleo{h_n}$, e.g.
  \begin{align}
    \seq{\epsilon_n} \in \littleo{h} \iff \lim_{n\to\infty} \abs{\frac{\epsilon_n}{h}} = \lim_{n\to\infty} \abs{\epsilon_n \sqrt{npq}} = 0
  \end{align}
\end{remark}

\section{Arithmetic rules}

All the presented notations share a common property which is in the case of the big O notation \cite{tampis:general}:
\begin{align}
  \seq{\epsilon_n} \in \bigo{a_n} \iff \seq{\abs{\frac{\epsilon_n}{a_n}}} \in \bigo{1}
\end{align}
This property of $\bigo{\cdot}$ can be condensed to the equation $\bigo{a_n} = a_n \cdot \bigo{1}$ because
\begin{align}
  \seq{\epsilon_n} \in \bigo{a_n} & \iff \seq{\abs{\frac{\epsilon_n}{a_n}}} \in a_n \cdot \bigo{1} \bnl
  & \iff \exists \seq{\delta_n} \in \bigo{1} : \abs{\frac{\epsilon_n}{a_n}} = \delta_n \bnl
  & \iff \exists \seq{\delta_n} \in \bigo{1}: \abs{\epsilon_n} = \abs{a_n} \delta_n \bnl
  &
  \begin{comment}
    \seq{a_n} \in \bigo{b_n} \implies \forall \seq{c_n} \in \{0,1\}^{\N} : \seq{(-1)^{c_n} a_n} \in \bigo{b_n}
  \end{comment} \bnl
  & \iff \exists \seq{\delta_n} \in \bigo{1} : \epsilon_n = a_n \cdot \delta_n \bnl
  & \iff \seq{\epsilon_n} \in a_n \cdot \bigo{1} \bnl
\end{align}

So $\bigo{a_n}$ is fully determined by knowing $\bigo{1}$ (cf. \cite{tampis:char}). For $\bigabs{\cdot}$ and $\littleo{\cdot}$ we get the analogous equation. Because $A(1)$ for $A\in\{O, o, \Psi, \Delta\}$ determines $A(a_n)$ we can deduce all arithmetic rules for asymptotic notation from the arithmetic rules for sets of the form $A(1)$. For example there are the rules \cite{tampis:rules}

\begin{enumerate}
  \item $(1)_{n\in\mathbb N} \in A(1) \implies (a_n)_{n\in\mathbb N} \in A(a_n)$
  \item $A(1) \subseteq B(1) \implies A(a_n) \subseteq B(a_n)$
  \item $A(1)\cdot B(1) \subseteq C(1) \implies A(a_n)\cdot B(b_n) \subseteq C(a_n\cdot b_n)$ \label{rule:example}
  \item $A(1)\cdot B(1) \subseteq C(1) \implies A(B(a_n)) \subseteq C(a_n)$ \label{rule:toproof}
  \item $A(1)+ B(1) \subseteq C(1) \implies A(a_n) + B(a_n) \subseteq C(a_n)$ \label{rule:demo}
\end{enumerate}

These rules can be mostly proven by using the equation $A(a_n) = a_n \cdot A(1)$. For example when we assume the premise $A(1) + B(1)\subseteq C(1)$ we can demonstrate rule \ref{rule:demo} via

\begin{align}
  A(a_n) + B(a_n) = a_n (A(1)+B(1)) \subseteq a_n C(1) = C(a_n)
\end{align}
Only the rule \ref{rule:toproof} is not obvious. However we find by assuming $A(1) \cdot B(1) \subseteq C(1)$ that
\begin{align}
  \seq{\epsilon_n} \in A(\delta_n) \land \seq{\delta_n} \in B(a_n) & \implies \seq{\abs{\frac{\epsilon_n}{\delta_n}}} \in A(1) \land \seq{\abs{\frac{\delta_n}{a_n}}} \in B(1) \bnl
  &
  \begin{comment}
    A(1) \cdot B(1) \subseteq C(1)
  \end{comment} \bnl
  & \implies \seq{\abs{\frac{\epsilon_n}{\delta_n}} \cdot \abs{\frac{\delta_n}{a_n}}} = \seq{\abs{\frac{\epsilon_n}{a_n}}} \in C(1) \bnl
  & \implies \seq{\epsilon_n} \in C(a_n)
\end{align}
and thus $A(B(a_n))\subseteq C(a_n)$. The arithmetic rules for sets $A(1)$ are mostly easy to prove or are already known theorems. For example $\littleo{1}\cdot\bigabs{1} \subseteq \littleo{1}$ is the well known proposition
\begin{align}
  \lim_{n\to\infty} \epsilon_n = 0 \land \sup_{n\in\N} \abs{\delta_n} \le 1 \implies \lim_{n\to\infty} \epsilon_n \delta_n = 0
\end{align}

From rule \ref{rule:example} and \ref{rule:toproof} we can now conclude that $\littleo{a_n} \cdot \bigabs{b_n} \subseteq \littleo{a_n \cdot b_n}$ and $\littleo{\bigabs{a_n}} \subseteq \littleo{a_n}$. We see that most of the arithmetic rules for asymptotic notations are corollaries from already known theorems about sequences and their limit behavior.

\section{Taylor series}

\begin{theorem}
  Let $f:D\to\R$ be a two times continuously differentiable function with an open domain $D$. Let $x_0\in D$. Then

  \begin{align}
    f(x_0+h) = f(x_0) + f'(x_0) h + \bigo{h^2}
  \end{align}
\end{theorem}

\begin{proof}
  By using the Lagrange form of the remainder we get the following Taylor expansion \cite[p. 284]{koenigsberger}
  \begin{align}
    f(x_0+h)=f(x_0)+f'(x_0)h + \frac 12 f''(\tilde x) h^2
  \end{align}
  with $\tilde x$ being in the interval bounded by $x_0$ and $x_0+h$. We fix a connected and compact neighborhood $U$ of $x_0$. $f''$ being continuous implies that $f''$ is bounded on $U$. Thus we find a $C > 0$ such that for all $h$ with $x_0+h\in U$ we find $\abs{f''(\tilde x)} \le C$. Hence for all $h$ with $x_0+h\in U$:
  \begin{align}
    f(x_0+h)& =f(x_0)+f'(x_0)h + \frac 12 f''(\tilde x) h^2 \nl
    & \in f(x_0)+f'(x_0)h + \bigabs{\frac 12 C h^2} \nl
    & \subseteq f(x_0)+f'(x_0)h + \bigo{h^2}
  \end{align}
\end{proof}

\section{Asymptotic behavior of Riemann sums}

Because we will have to deal with Riemann sums we need the following theorem:

\begin{theorem} \label{thm:asymptotics_riemann}
  Let $h > 0$ and $\x[k]$ be defined via $\x[k] = \x[0] + kh$ for $k\in\Z$ and any $\x[0]\in\R$. Let $f:\R\to\R$ with $\snorm{f} < \infty$ and $\snorm{f'} < \infty$. Let $a,b\in\R$. We find

  \begin{enumerate}
    \item $\int_x^{x+h} f(t) \d{t} \in f(x)h + \bigabs{\frac 12 \snorm{f'} h^2}$ and $\int_{x-h}^x f(t) \d{t} \in f(x) h + \bigabs{\frac 12 \snorm{f'} h^2}$
    \item $\f{\x} h \in \int_{\x-\frac h2}^{\x + \frac h2} \f{t}\d{t} + \bigabs{\frac 14 \snorm{f'} h^2}$
    \item For $a \le b$ we have

      \begin{align}
        \sum_{a \le \x \le b} \f{\x} h \in \int_a^b \f{t}\d{t} + \bigabs{\snorm{f} h + \frac 14 \snorm{f'}(b-a)h + \frac 14 \snorm{f'}h^2}
      \end{align}
  \end{enumerate}
\end{theorem}

\begin{proof}~
  \begin{enumerate}
    \item We find
      \begin{align}
        \abs{\int_x^{x+h} \f{t} \d{t} - f(x)h} &= \abs{\int_x^{x+h} \br{\f{t}-\f{x}} \d{t} } \bnl
        &
        \begin{comment}
          \forall t\,\exists \epsilon_t: \f{t}-\f{x}=\fp{\epsilon_t}(t-x)
        \end{comment} \bnl
        & \le \int_{x}^{x+h} \abs{\fp{\epsilon_t}} \abs{t-x} \d{t} \bnl
        & \le \snorm{f'} \int_{0}^{h} \abs{t} \d{t} \bnl
        & = \frac 12 \snorm{f'} h^2
      \end{align}
      Analogously we can demonstrate that $\int_{x-h}^x f(t) \d{t} \in f(x) h + \bigabs{\frac 12 \snorm{f'} h^2}$.

    \item We have
      \begin{align}
        \f{\x} h & = \f{\x} \frac h2 + \f{\x} \frac h2 \bnl
        &\in \int_{\x}^{\x+\frac h2} \f{t}\d{t} + \bigabs{\frac 18 \snorm{f'} h^2} + \int_{\x-\frac h2}^{\x} \f{t}\d{t} + \bigabs{\frac 18 \snorm{f'} h^2} \bnl
        &\subseteq \int_{\x-\frac h2}^{\x+\frac h2} \f{t}\d{t} + \bigabs{\frac 14 \snorm{f'} h^2}
      \end{align}

    \item Let $\tilde a = \inf \{ \x : \x \ge a \}$ and $\tilde b = \sup \{ \x : \x \le b \}$. The difference of $a$ and $\tilde a$ is less than or equal $h$ and the distance of $a$ and $\tilde a - \frac h2$ can be estimated by $\frac h2$. Similarly we find $\abs{\tilde b-b} \le h$ and $\abs{\br{\tilde b+\frac h2}-b} \le \frac h2$. We have

      \begin{align}
        \sum_{a \le \x \le b} \f{\x} h & \in \sum_{a\le \x\le b} \br{\int_{\x-\frac h2}^{\x+\frac h2} \f{t} \d{t} + \bigabs{\frac 14 \snorm{f'} h^2}} \bnl
        & \subseteq \int_{\tilde a - \frac h2}^{\tilde b + \frac h2} \f{t} \d{t} + \bigabs{\frac 14 \snorm{f'}\br{\tilde b + h - \tilde a} h}\bnl
        & \subseteq \int_{a+\bigabs{\frac h2}}^{b+\bigabs{\frac h2}} \f{t} \d{t} + \bigabs{\frac 14\snorm{f'}\br{b-a+\bigabs{h}} h}\bnl
        & \subseteq \int_{a+\bigabs{\frac h2}}^{b+\bigabs{\frac h2}} \f{t} \d{t} + \bigabs{ \frac 14 \snorm{f'} (b-a) h + \frac 14 \snorm{f'} h^2} \bnl
        &
        \begin{comment}
          \begin{aligned}
            \int_a^{b+\bigabs{\epsilon}} \f{t} \d{t} & \subseteq \int_a^b \f{t} \d{t} + \int_b^{b+\bigabs{\epsilon}} \f{t} \d{t} \bnl
            & \subseteq \int_a^b \f{t} \d{t} + \snorm{f} \int_b^{b+\bigabs{\epsilon}} \d{t} \bnl
            & \subseteq \int_a^b \f{t} \d{t} + \snorm{f} \bigabs{\epsilon}
          \end{aligned}
        \end{comment} \bnl
        & \subseteq \int_a^b \f{t} \d{t} + \bigabs{\snorm{f}h + \frac 14 \snorm{f'} (b-a) h + \frac 14 \snorm{f'} h^2}
      \end{align}
  \end{enumerate}
\end{proof}


  \chapter{The binomial distribution}

\section{The motivational problem for the binomial distribution}

Probability theory emerged from and was highly influenced by investigations of games of chance~\cite[p. 4]{hald1}. Since the antiquity dice games were played and starting from the 14th century card games became more and more popular~\cite[pp. 33-34]{hald1}. Governments used lotteries to finance their expenditures and a lot of private lotteries were conducted as well~\cite[p. 34]{hald1}.

\includewrapfig{Christiaan_Huygens}{Christiaan Huygens}{File \href{https://commons.wikimedia.org/wiki/File:Christiaan_Huygens.jpg}{``Christiaan Huygens.jpg''} from Wikimedia Commons uploaded by \href{https://commons.wikimedia.org/w/index.php?title=User:Lord_Horatio_Nelson~commonswiki}{Lord Horatio Nelson$\sim$commonswiki} and licensed under Public domain}

From the economic and recreational importance arose a demand in calculating the odds of a game or the value of the expected winnings~\cite[p. 34]{hald1}. In 1654 Fermat and Pascal solved in a correspondence the problem of division~\cite[pp. 42-64]{hald1} which marks the foundation of probability theory~\cite[p. 4]{hald1}\footnote{In the 16th century Cardano already discussed several problems about games of chance in his work \emph{Liber de ludo aleae} (Book on Games of Chance)~\todo{cite, name des Buchs}\cite[pp. 33-41]{hald1}. However his book was first published posthumously 1663 and thus did not influenced Fermat, Pascal or Huygens~\cite[p. vii]{bernoulli}.}. Christiaan Huygens, who (according to himself) heard from the letters by Pascal and Fermat but did not know their methods~\cite[p. 67]{hald1}, wrote a short treatise \emph{De Ratiociniis in Ludo Aleae} (On Reckoning in Games of Chance) which was published 1657~\cite[p. vii]{bernoulli}\footnote{Huygens wrote his treatise in Dutch. His mathematics teacher Frans van Schooten translated it in Latin and published it at the end of his book \emph{Exercitationes Mathematicae}~\cite[pp. 65-68]{hald1}.}. Among the different problems, Huygens solved in his work, was the following~\cite[p. 163]{bernoulli}:

\begin{quotation}
  To find with how many dice one may undertake to throw two sixes on the first try.
\end{quotation}

\includewrapfig[0.5\textwidth]{tavern1658}{Men gambling in a tavern (image section of a Flemish painting from 1658)}{Painting \href{https://commons.wikimedia.org/wiki/File:Tavern_Scene-1658-David_Teniers_II.jpg}{``Tavern Scene''} by \href{https://en.wikipedia.org/wiki/David_Teniers_the_Younger}{David Teniers the Younger} from 1658. Cropped version of picture which was uploaded at Wikimedia Commons and which was taken by National Gallery of Art, Washington, D.C., USA. Picture is licensed under Public Domain.}

The problem is to calculate the number of dices one need to throw such that the probability of throwing two sixes is at least $\tfrac 12$. Imagine for example that you are in a tavern of the 17th century drinking beer. A merchant comes to your table and offers you the following deal: ``I'll give you $8$ dices which you can throw at once. If you will get at least two sixes you will get a gold coin. In case you only throw one six or none you have to pay me a gold coin.'' Shall you participate in his game? Is the game fair or not? Who has the higher chance to win? This example demonstrates very well the motivation of Huygens and his colleagues to engage in such problems. 

To solve this problem one first need to calculate the probability two throw exactly $m$ sixes with $n$ dices. If you are already familiar with the binomial distribution you see how this distribution can be used in the solution of the problem. This was later done by Jacob Bernoulli who reprinted Huygens' work in the first part of his book \emph{Ars Conjectandi} (The Art of Conjecturing) where he also added additional comments~\cite[p. 63]{bernoulli}. Before we will have a look at Bernoulli's solution with a derivation of the binomial distribution I want to show how Huygens dealt with the problem to have a comparison of both attempts. Huygens first noted~\cite[p. 163]{bernoulli}:

\begin{quotation}
  Now this is just the same as asking in how many throws a person may undertake to throw one die in order to get two sixes.
\end{quotation}

So it does not make any difference whether $n$ dices are thrown at once or whether there is one die which is thrown $n$ times. Jacob Bernoulli gave a good explanation for this circumstance in his reprint of Huygens' work~\cite[p. 163]{bernoulli}:

\begin{quotation}
  If, for example, one throw of ten dice is allowed, then it is certainly evident that it makes no difference whether those ten dice are thrown onto the gaming board altogether at one time or successively one after another. And if it is done successively, then it is equally clear that it makes no difference whether the ten dice thrown are ten different dice or one and the same die retrieved from the board and thrown ten times.
\end{quotation}

After noting that it does not matter how the dices are thrown, Huygens showed how to calculate the expected winning recursively when one wins the stake $a$ for throwing at least two sixes in a certain number of throws~\cite[p. 163]{bernoulli}:

\begin{quotation}
  If someone undertook to do this in two throws, there would fall to him $a/36$, by what was shown before. If he is given three throws, then, if his first throw is a six, he will still have two throws, both of which must be sixes, and we said this is worth the same as $a/36$. But if his first throw is not a six, then he needs only to get one six on the remaining two throws. By proposition X this is worth of having $11a/36$. But there is certainly one case in which he throws a six th first time and five cases in which something else happens. So in the beginning there is one case for $11a/36$ and five cases for $a/36$. By proposition III\todo{typo untersuchen}, this is worth as much as $16a/127$ or $2a/27$. In repeatedly considering one more throw, we find that there is an advantage in undertaking to throw two sixes in ten throws with one die or in one throw with ten dice.
\end{quotation}

The argumentation of Huygens is the following: First he looked at the case when two dices a thrown. In proposition XI he already had calculated the expected value for this case to be $\tfrac 1{36}a$. The he investigated the case of three throws. There are five cases in which the first throw is not a six and the remaining expected value in these cases is $\tfrac 1{36}a$ because two sixes must be thrown in the remaining two throws. In one case the first throw is a six. Huygens already had shown in proposition X that the expected value of throwing at least one six in two throws is $\tfrac{11}{36}a$. So in total there a $5$ cases with a remaining expected value of $\tfrac 1{36}a$ and $1$ case with the remaining value of $\tfrac{11}{36} a$. The total expected value for three throws was now calculated as

\begin{align}
  \frac{1}{5+1} \left(5 \cdot \frac 1{36} a + 1 \cdot \frac{11}{36} a \right) = \frac{1}{6} \cdot \frac{5+11}{36} a = \frac{16}{216} a = \frac{2}{27} a
\end{align}

\noindent Thereby Huygens used his proposition III which states~\cite[p. 135]{bernoulli}:

\begin{quotation}
  If the number of cases in which $a$ falls to me is $p$ and the number of cases in which $b$ falls to me is $q$, and if all the cases can happen equally easily, then my expectation will be worth $(pa+qb)/(p+q)$.\footnote{The final formula can be rewritten as $\tfrac{p}{p+q} a+\tfrac q{p+q} b$ which is the definition of the expected value for the Bernoulli trial, a distribution with two possible outcomes. \todo{verbessern}}
\end{quotation}

Without writing it explicitly down in his book, Huygens now calculates recursively the odds of having two sixes in $n$ throws. His calculation might have looked like the following: Let $O_n$ be the expected winning of having at least one six in $n$ throws. In proposition X Huygens already had demonstrated by looking at several examples that $O_n = \tfrac 1{1+5} \left( 1 \cdot a + 5 \cdot O_{n-1} \right)=\tfrac 16 \left( a + 5 \cdot O_{n-1} \right)$. If the first throw is a six, then the player gets the whole stack $a$. In the remaining five cases the player still has the opportunity to throw a six in the remaining $n-1$ throws and thus have an expected winning of $O_{n-1}$ to win. The calculated values by Huygens are:

\begin{align}
  \begin{array}{rll}
    O_1 & = \frac 16 a \nl
    O_2 & = \frac 16 ( a + 5 \cdot \frac 16 a) & = \frac{11}{36} a \nl
    O_3 & = \frac 16 ( a + 5 \cdot \frac{11}{36} a) & = \frac{91}{216} a \nl
    O_4 & = \frac 16 ( a + 5 \cdot \frac{91}{216} a) & = \frac{671}{1296} a \nl
    & \vdots \nl
    O_8 & = \frac 16 ( a + 5 \cdot \frac{201811}{279936} a) & = \frac{1288991}{1679616} a \nl
    O_9 & = \frac 16 ( a + 5 \cdot \frac{1288991}{1679616} a) & = \frac{8124571}{10077696} a \nl
    & \vdots
  \end{array}
\end{align}

Let $T_n$ be the expected winning for throwing at least two sixes in $n$ throws. By the same argument Huygens gave for the special case $n=3$ we deduce $T_n = \tfrac{1}{6} ( 5 \cdot T_{n-1}  + O_{n-1})$. Thus Huygens calculated

\begin{align}
  \begin{array}{rlll}
    T_2 & = \frac 1{36} a \nl
    T_3 & = \frac 16 (5 \cdot \frac 1{36} a + \frac{11}{36} a) &= \frac 2{27} a \nl
    T_4 & = \frac 16 (5 \cdot \frac 2{27} a + \frac{91}{216} a) &= \frac {19}{144} a \nl
    & \vdots \nl
    T_9 & = \frac 16 (5 \cdot \frac {663991}{1679616} a + \frac{1288991}{1679616} a) &= \frac {2304473}{5038848} a & \approx 0.457 a \nl
    T_{10} & = \frac 16 (5 \cdot \frac{2304473}{5038848} a + \frac{8124571}{10077696} a) &= \frac{10389767}{20155392} a & \approx 0.515 a \nl
    & \vdots 
  \end{array}
\end{align}

So starting with $10$ dices the odds for the player to win is higher than $\tfrac 12$. Besides the probability for us to win in the above supposed game with eight dices is $\tfrac{663991}{1679616}$ which is roughly $0.40$.

\section{Bernoulli's derivation of the binomial distribution}

\includewrapfig{Jakob_Bernoulli}{Jacob Bernoulli}{File \href{https://commons.wikimedia.org/wiki/File:Jakob_Bernoulli.jpg}{``Jakob Bernoulli.jpg''} from Wikimedia Commons uploaded by user \href{https://commons.wikimedia.org/wiki/User:Materialscientist}{Materialscientist} and licensed under Public Domain}

As we have seen in the previous section Christiaan Huygens could well solve his proposed problem. One also can extract an algorithm from his method with which similar problems can be solved. But he only investigated a special case although a more general formulation of his method is possible. For example one may be interested in a game with dices which have more or less than six faces. Also the number of favorable faces on a die may be more then one and the minimal value of favorable results in all throws may differ from two. To the previous and similar problem Jacob Bernoulli wrote~\cite[p. 157]{bernoulli}:

\begin{quotation}
  If the Author [Christiaan Huygens, S.K.]\todo{richtig?} had substituted letters for numbers, he could have expressed this and the preceding proposition as one problem and investigated its general solution with equal ease [...]\todo{richtig?}
\end{quotation}

Also Huygens' recursive method needs a lot of calculational effort to find the final solution. Thus one question arises: Is there a more direct and general solution for Huygens' problem?

\includefig{Ars_Conjectandi}{Cover of \emph{Ars Conjectandi}}{Cover of \emph{Ars Conjectandi} by Jacob Bernoulli 1713. File \href{https://commons.wikimedia.org/wiki/File:Arsconj.gif}{``Arsconj.gif''} from Wikimedia Commons uploaded by user \href{https://commons.wikimedia.org/wiki/User:Nousernamesleft}{Nousernamesleft} and licensed under Public Domain.}

Jacob Bernoulli addressed this question with the binomial distribution as an answer in his book \emph{Ars Conjectandi} (The Art of Conjecturing) from 1713~\cite[pp. 220-256]{hald1}\footnote{The book was published eight years after Jacob Bernoulli's death by his nephew Nicholas Bernoulli in 1703~\cite[pp. 223-224]{hald1}}. In the first part of this book Bernoulli\footnote{The family of Bernoullis is huge with a lot of influential mathematicians. In the following I always refer to Jacob Bernoulli when I only mention ``Bernoulli''. A good overview of Bernoulli's family tree gives~\cite[pp. 1-4]{bernoulli}} reprinted Huygens' \emph{De Ratiociniis in Ludo Aleae} with additional notes and alternative solutions~\cite[p. 63]{bernoulli}.

In his comments to the above problem Bernoulli first derived the binomial distribution by incomplete induction with an argumentation similar to Huygens' solution. Afterwards he derived the binomial distribution with a combinatorial argumentation~\cite[pp. 165-167]{bernoulli}:

\begin{quotation}
  [Consider] $n$ dice A, B, C, D, etc., each constructed with $a$ faces, $b$ of which match the purpose of the player, and the other $c$ of which do not. And ask in how many cases it may happen that what was undertaken is accomplished on none of the dice, on just one of the dice, or on only 2, 3, or 4 of them, and so on; finally, in how many cases the person undertaking the game fails in his purpose and his opponent wins. But it was shown in the Note to the preceding Proposition that there are $C^n$ cases in which what was undertaken is accomplished on none of the dice. And it may be concluded in a similar way that there are $b$ cases, or $bb$, or $b^3$, etc., in which one of the dice, for instance A, or two, A and B, or three, A, B, and C, and so on, meet the purpose of the person undertaking the game. Similarly, there are $c^{n-1}$ cases, or $c^{n-2}$ cases, and so on, in which $n-1$ dice, or $n-2$ dice, or $n-3$ dice, and so on, fail to fulfill his hopes. Therefore, since each of these cases can be joined with any of the preceding, multiplication of the latter by the former produces $bc^{n-1}$ cases, or $bbc{n-2}$ cases, or $b^3c^{n-3}$ cases, and so on. And since the die or dice that favor the person undertaking the game might be either A or B or C, and so on, if there is one, or, if there are two, either A and B, or A and C, or B and C, and so on, or, if there are three, either A, B, and C, or A, B and D, and so on, these number of cases must again be multiplied by the number of units, pairs, triples, and so son that can be drawn from all $n$ dice. But by the theory of combinations developed in Part II, this number is $n$, or $n(n-1)/(1\cdot 2)$, or $n(n-1)(n-2)/(1\cdot 2\cdot 3)$, etc. So from this multiplication we obtain $nbc^{n-1}$, or

  
  $\frac{n(n-1)}{1\cdot 2}bbc^{n-2}$,

  or $\frac{n(n-1)(n-2)}{1\cdot 2\cdot 3}b^3c^{n-3}$ etc., on up to

  $\frac{n(n-1)(n-2)\ldots(n-m+2)}{1\cdot 2\cdot 3\cdot 4\ldots(m-1)} b^{m-1} c^{n-m+1}$

  for the numbers of cases in which what was undertaken is accomplished with exactly one die, or two dice, or three dice and so on up to $m-1$ dice. Since, as we mentioned, the opponent wins in all these cases and since there are altogether $a^n$ cases with $n$ dice, his lot is, \textbf{[42]} by Corollary 1 of Proposition III,

  $c^n+nbc^{n-1}+\frac{n(n-1)}{1\cdot 2}bbc^{n-2} + \frac{n(n-1)(n-2)}{1\cdot 2\cdot 3}b^3c^{n-3} \ldots$

  $+\frac{n(n-1)(n-2)\ldots(n-m+2)}{1\cdot 2\cdot 3\cdot 4\ldots(m-1)} b^{m-1}c^{n-m+1} : a^n$,\footnote{The notation $x_1+x_2+\ldots+x_n:y$ means $\tfrac 1y (x_1+x_2+\ldots+x_n)$. \todo{cite}}

  as above.
\end{quotation}

Bernoulli considered a die with $a$ faces which has $b$ faces being favorable for the player and $c=a-b$ faces being unfavorable. The player has to throw at least $m$ times one of the $b$ favorable faces in a total of $n$ throws in order to win. What is the probability of winning / loosing?

Bernoulli started with calculating the number of outcomes in which the player throws exactly $k$ times one of the favorable faces. Fix $k$ of the total $n$ throws. All of these $k$ throws shall be one of the $b$ favorable results and the other $n-k$ shall be one of the $c$ unfavorable results. Because each of these $k$ throws can show one of $b$ possible faces the total number of different outcomes in these $k$ throws with only favorable results is $b^k$. Analogously is the number of possible outcomes in the other $n-k$ throws with only unfavorable results equal to $c^{n-k}$. Therefore there are in total $b^kc^{n-k}$ different outcomes in which the $k$ fixed throws show a favorable face and the other $n-k$ throws show a unfavorable face.

In favor of throwing exactly $k$ favorable numbers, the order in which those $k$ numbers are thrown does not matter. As known by combinatorics there are $\binom nk=\frac{n!}{k!(n-k)!}=\frac{n\cdot(n-1)\ldots(n-k+1)}{1\cdot2\cdot3\ldots k}$ different orders in which exactly $k$ favorable throws occur. For each of these orders the total number of different outcomes, in which the $k$ favorable and $n-k$ unfavorable results occur in the specified order, is $b^kc^{n-k}$ as we have seen in the previous paragraph. Thus the total numbers of outcomes with exactly $k$ favorable throws is $\binom nk b^k c^{n-k}$.

Because the total number of outcomes is $a^n$ and because all outcomes are equally likely, the probability for the player to have exactly $k$ favorable throws is $\frac {1}{a^n} \binom nk b^k c^{n-k}$.

Because the player looses when he throws at most $m-1$ favorable numbers, the probability for the player to loose is

\begin{align}
  \P{\text{player looses}} = \frac{1}{a^n} \sum_{k=0}^{m-1} \binom nk b^k c^{n-k}
\end{align}

This is the same formula Bernoulli derived in the end of the above quoted section. He ended his discussion of Huygens' problem with two corollaries which he named ``Rule for knowing the lot of a player to whom several throws of a die have given and who is held to archieving something precisely on some certain throws and not on others''~\cite[p. 170]{bernoulli}:

\begin{quotation}
  \emph{Corollary 1.} If there are the same numbers of cases in all games [...] the expectation found to be [...] $b^mc^{n-m}/a^n$, taking $n$ for the number of all the throws, and $m$ for the number in which what is to be achieved ought to be achieved.
  
  \emph{Corollary 2.} If the same numbers of cases rule in all games and if, in addition, the number of games or throws is determined within which something is to be achieved, but such that the desired result can occur on any throws and not on certain determined ones [...] then it is clear that the quantity of the expectation found in the previous corollary should be multiplied by as many times as [...] $m$ things may be chosen from $n$ diverse things. This, however, may occur, by the Theory of Combinatorics to be treated in Part II, in $n(n-1)(n-2)(n-3)\ldots(n-m+1)/(1\cdot 2\cdot 3\cdot 4\ldots m)$ ways [...]. Thus the expectation of the player will be worth

  \begin{align*}
    \frac{n(n-1)(n-2)(n-3)\ldots(n-m+1)}{(1\cdot 2\cdot 3\cdot 4\ldots m)}\cdot \frac{b^mc^{n-m}}{a^n}
  \end{align*}
\end{quotation}

When we introduce $p=\tfrac ba$ which is the probability that one get a favorable result in one throw and analogously $q=\tfrac ca=1-p$ for the probability of an unfavorable result in a throw, we get:

\begin{align}
  \P{\text{$m$ in $n$ throws are favorable}} & = \frac{n(n-1)(n-2)(n-3)\ldots(n-m+1)}{(1\cdot 2\cdot 3\cdot 4\ldots m)}\cdot \frac{b^mc^{n-m}}{a^n} \nl
  &= \binom nm \frac{b^m}{a^m} \frac{c^{n-m}}{a^{n-m}} \nl
  &= \binom nm p^m q^{n-m}
\end{align}

So we finally get a formula for calculating the odds of having $m$ favorable cases in a total of $n$ (independent) cases when in each case $p$ is the probability of having a favorable case\footnote{In Bernoulli's work $p$ is $\tfrac ba$ and thus Bernoulli only considered rational probabilities in his derivation. This was common in the beginning of probability theory \todo{cite}. It may have its origin, that in most games with a finite number of rounds only rational probabilities occur.}. As we will see in the next section, this formula gets the probability mass of the binomial distribution.

\section{Definition of the binomial distribution}

In the previous section we have seen the historical derivation and motivation of the binomial distribution. In this section I want to introduce the binomial distribution in a way contemporary textbooks define it. \todo{cite}

Take $n$ independent random variables $X_1, X_2, \ldots, X_n$. Each of them has two possible outcomes $\alpha$ and $\beta$. Those random variables with two possible outcomes are called ``Bernoulli trials'' named after Jacob Bernoulli. \todo{cite} The outcome $\alpha$ has the probability $p$ and $\beta$ has the probability $q=1-p$ to occur in each random variable $X_i$. 

The binomial distribution is the probability distribution which assigns for each $k\in \{0,1,\ldots,n\}$ the probability of having exactly $k$ times the outcome $\alpha$ among all random variables $X_1,X_2,\ldots,X_n$. A typical way to define the binomial distribution is to set $\alpha=1$ and $\beta=0$. The random variable $\Bs=\sum_{i=1}^n X_i$ is then binomially distributed.

What is $\P{\Bs = k} = \P{\sum_{i=1}^n X_i=k}$? To find this probability we can argue in a very similar way as in the above section. First we fix an order in which $k$ random variables of all $X_i$ are $\alpha$ and the rest $n-k$ random variables have $\beta$ as an outcome. The probability for each such order is $p^kq^{n-k}$ since the random variables are independent.

For example the first $k$ random variables $X_1,X_2,\ldots,X_k$ might have the result $\alpha$ and the other $n-k$ random variables $X_{k+1},X_{k+2},\ldots,X_n$ have the outcome $\beta$. The probability in this case is:

\begin{align}
  & \P{X_1=\alpha, X_2=\alpha, \ldots, X_k=\alpha,X_{k+1}=\beta, X_{k+2}=\beta, \ldots X_n =\beta} \nl
  =\ & \underbrace{\P{X_1=\alpha} \cdot \P{ X_2=\alpha} \ldots \P{X_k=\alpha}}_{=p^k} \cdot \underbrace{\P{X_{k+1}=\beta} \cdot \P{X_{k+2}=\beta} \ldots \P{X_n =\beta}}_{= q^{n-k}} \nl
  =\ & p^kq^{n-k}
\end{align}

Since the order in which the $k$ outcomes $\alpha$ and the $n-k$ outcomes $\beta$ occur does not matter, we have to sum of the probabilities over all such orders. Do define such an order, we have to take $k$ among all $n$ random variables which shall have the result $\alpha$. From combinatorics we know that there are $\binom nk$ ways to take $k$ objects from a whole set of $n$ objects. So there are $\binom bk$ orders in which exactly $k$ random variables have the outcome $\alpha$. Each of these orders has the possibility $p^kq^{n-k}$. So the total possibility to have exactly $k$ times the outcome $\alpha$ among all random variables is $\binom nk p^kq^{n-k}$.

As already mentioned, this distribution is called ``binomial distribution'':

\begin{definition}[Binomial distribution]
  The binomial distribution $\Bs$ is the probability distribution on the set $\{0,1,2,\ldots,n\}$ with the probability mass function

  \begin{align}
    \P{\Bs=k} = \binom nk p^k q^{n-k}
  \end{align}
\end{definition}

\includefig{binomial_distribution_pmf}{The probability mass function of the binomial distribution}{File \href{https://commons.wikimedia.org/wiki/File:Binomial_distribution_pmf.svg}{``Binomial distribution pmf.svg''} from Wikimedia Commons. Upload by user \href{https://commons.wikimedia.org/w/index.php?title=User:Tayste}{Tayste} and licensed under Public Domain.}

\section{Applications of the binomial distribution}

The binomial distribution occurs often in statistics and thus in all fields which uses statistical tools or probability theory like biology~\cite[p. 146]{henze}, psychology, economics, physics etc. \todo{cite}. In the beginning section of this chapter we have already met an application of this distribution in gambling too.

\includewrapfig{assembly_line}{An assembly line}{File \href{https://www.flickr.com/photos/32659528@N00/2868489384}{``Final assembly''} from Flickr by \href{http://flickr.com/photos/32659528@N00}{Brian Snelson} and licensed under \href{https://creativecommons.org/licenses/by/2.0/deed.de}{CC-BY 2.0}.}

The binomial distribution occurs often in counting processes~\cite[p. 138]{henze} to calculate the odds of having $k$ times a certain property among $n$ objects or occasions. The binomial distribution is also used as a simplified model when the probability of having a certain property varies only a little bit between all objects/occasions or when there is only a slight dependence in the occurence of the property between all objects/occasions. Take the following example:

\begin{example}
  What is the probability that $k$ among $n$ made products in an assembly line are faulty if the probability of one product produced incorrectly is $p$.
\end{example}

Firstly it is likely that the probability of having a faulty product varies over time. For example it may depend on the temperature of the machines which may fluctuate. Secondly the occurrence of a badly made product may increase the likelihood of flawed products for the next produced merchandises. There may be a failure in one of the machines which also effects the next products until the failure is corrected.

Although the premises for the binomial distribution is not fulfilled in this example, this distribution may be used by the company to do statistical analysis. This shows how important the binomial distribution for contemporary statistical analysis is because it may be used whenever one investigate a counting process.

\section{The binomial distribution and the sex ratio}

\includewrapfig{great_plague}{The great plague in London, 1665}{File \href{https://commons.wikimedia.org/wiki/File:Great_plague_of_london-1665.jpg}{``Great plague of london-1665.jpg''} from Wikimedia Commons uploaded by user \href{https://commons.wikimedia.org/wiki/User:7mike5000}{7mike5000} and licensed under Public Domain.}

I want to finish this chapter with presenting you how early statisticians used the binomial distribution to investigate the sex ratio at birth. This investigation was not only one of the earliest statistical analysis it also motivates the theorem by De Moivre and Laplace.

Starting from the 14th century plague epidemics occurred frequently in Europe for nearly 400 years~\cite[p. 82]{hald1}. So since the 1530s the parish clerks in London were requested to publish weekly the number of plague deaths together with the number of all other deaths~\cite[p. 82]{hald1}. These numbers were used by the government to recognize a new plague epidemic so that countermeasures could be taken. Also the wealthier part of the society used this data as an indication when they shall leave London~\cite[pp. 82-83]{hald1}. Alongside the number of deaths, the number of weddings and christenings inside the Church of England were registered too~\cite[p. 83]{hald1}. Besides the burials and christenings were given separately for male and female starting from 1629~\cite[p. 83]{hald1}.

The first statistical analysis of this data were made by John Graunt~\cite[p. 83]{hald1} which was also one of the first statistical analysis of history~\cite[p. 81]{hald1}. In 1662 he published his book \emph{Natural and Political Observations Made upon the Bills of Mortality} in which he summarized and investigated the published data about mortality and birth in London~\cite[pp. 81-105]{hald1}. He calculates the sex ratio in the christenings for each $8$ years and concluded that there are in average more men than female born~\cite[pp. 92-93]{hald1}.

\includewrapfig{bill_of_mortality}{The bill of mortality for 1665}{File \href{https://commons.wikimedia.org/wiki/File:Bill_of_Mortality.jpg}{``Bill of Mortality.jpg''} from Wikimedia Commons. Uploaded by user \href{https://commons.wikimedia.org/wiki/User:Dolledre}{Dolledre} and licensed under Public Domain.}

\begin{figure}
  \begin{center}
    \begin{tabular}{c|c|c}
      Year & Male & Female \\
      \hline
      1629 & 5218 & 4683 \\
      1630 & 4858 & 4457 \\
      1631 & 4422 & 4102 \\
      1632 & 4994 & 4590 \\
      1633 & 5258 & 4839 \\
      \vdots & \vdots & \vdots
    \end{tabular}

    \caption{Samples of the used data by Arbuthnott~\cite[p. 276]{hald1}}
  \end{center}
\end{figure}

In the following many mathematicians were inspired by Graunt's book to describe the reported phenomenas with probabilistic models~\cite[p. 275]{hald1}. One of these mathematicians was John Arbuthnott. He summarized the yearly number of male and female birth for 82 years in the period 1629-1710 and observed that for each year there are more boys than girls born. In this disproportion he saw a sign of the Divine Design and an evidence of the existence of God~\cite[p. 277]{hald1}\footnote{Arbuthnott's belief was common in the beginning of the 18th century among statisticians and theologians~\cite[p.~285]{hald1}. This may have motivated investigations about the sex ratio during this time.}. He argued that the chance for a boy to die during growing up is slightly higher than for girls. So the bigger chance for a male birth compensates this higher mortality rate in childhood so that there is an equal number of males and females and adulthood~\cite[p. 277]{hald1}. John Arbuthnott wrote~\cite[p. 275]{hald1}:

\begin{quotation}
    This Equality of [adult] Males and Females is not the Effect of Chance but Divine Providence.
\end{quotation}

To prove his assumptions that the chances for male offspring is (a little bit) higher than for female offspring, Arbuthnott made the first test of significance in history~\cite[p. 276]{hald1}. He identified the determination of the birth's sex with the throw of a two sided die~\cite[pp. 275-276]{hald1}. Let the chances of a male birth be $p$ and thus let $q$ be the chances of the female birth\footnote{Here we ignore the possibility of an intersex birth. \todo{ausbauen}}. The yearly numbers of male and female births are then binomially distributed.

Arbuthnott then estimated the probability that the number of male births exceeded the number of female births for 82 consecutive years under the binomial model for $p=\tfrac 12$. He argued that because of the symmetry of the binomial distribution for $p=\tfrac 12$ the possibility for having more boys born than girls in one year is less or equal to $\tfrac 12$. Thus the possibility that this happens for $82$ consecutive years must be less than or equal $\left(\tfrac 12\right)^{82}$ which is less than $2.1\cdot 10^{-25}$. Because this possibility is so low the assumption $p=\tfrac 12$ must be rejected for the alternative $p > \tfrac 12$~\cite[p. 278]{hald1}.

  \chapter{The De Moivre-Laplace theorem}

\section{Formulation of De Moivre-Laplace theorem}

\includefig[0.25\textwidth]{moivre}{Abraham de Moivre}{File \href{https://commons.wikimedia.org/wiki/File:Abraham_de_moivre.jpg}{``Abraham de moivre.jpg''} from Wikimedia Commons. Uploaded by user \href{https://commons.wikimedia.org/wiki/User:\%E7\%AB\%B9\%E9\%BA\%A6\%E9\%AD\%9A(Searobin)} Searobin and licensed under Public Domain. Original source: \url{https://www.york.ac.uk/depts/maths/histstat/people/} .}

\includefig[0.25\textwidth]{laplace}{Pierre Simon Laplace}{File \href{http://flickr.com/photos/37667416@N04/4840056407}{``Laplace''} from Flickr. Uploaded by \href{http://www.flickr.com/people/37667416@N04}{Biblioteca de la Facultad de Derecho y Ciencias del Trabajo Universidad de Sevilla} and licensed under \href{https://creativecommons.org/licenses/by/2.0/deed.en}{CC-BY 2.0}.}

In the last chapter we have seen with the application of the binomial distribution in the investigation of the sex ratio that calculations with the binomial distribution are arduous. In order to compute $\binom nk p^kq^{n-k}$ there are $n+2\min\{k,n-k\}-1$ operations necessary.

This motivates approximations to the binomial distribution to minimize the computational efforts. With those approximations also equations like $\P{\Bs \le z}=\alpha$ are easier to solver for $z$ or $n$ ($\Bs$ shall be binomially distributed). Explicit solutions for those equations are hard to calculate~\cite[p. 469]{hald1}.

The first simple approximation to the binomial distribution was found by Abraham de Moivre with the help of James Stirling~\cite[p. 469]{hald1}. His proof was later extended by Pierre Simon Laplace~\cite[pp. 495 ff.]{hald1}. Today this approximation is known as the ``De Moivre-Laplace theorem''~\cite[pp. 64-67]{irle}. \todo{besseres Zitat}

I want to split this theorem in two parts. With the first part one get an approximation for the probability mass function of a binomial distribution:

\begin{theorem}[Local version of De Moivre-Laplace theorem]
  Let $p\in(0,1)$. ...
\end{theorem}

In the second part one gets an approximation to the culmulative distribution function which is the actual theorem by De Moivre and Laplace:

\begin{theorem}[De Moivre-Laplace theorem]
  ...
\end{theorem}

\section{The journey to De Moivre-Laplace theorem}

In \emph{A history of propability and statistics and their applications before 1750} Anders Hald gave a good summarize how de Moivre and Laplace found their theorem. In \cite[pp. 469-470]{hald1} he wrote:

\begin{quotation}
  In 1721 de Moivre began his investigations od the binomial distribution for $p=\tfrac 12$. He first found an approximation to the maximum term and next an approximation to the ratio of the maximum to the term at a distance of $d$ from the maximum. [...] From 1725 onward James Stirling (1692-1770) worked on the same problem and found that the constant entering de Moivre's formula equals $\sqrt{2\pi}$. After having obtained these results they realized that it would be simpler to begin with an approximation to $\ln n!$, and they both proved Stirling's formula, $n! \sim \sqrt{2\pi n}n^ne^{-n}$. De Moivre's proofs are given in the \emph{Miscellanea Analytica} (MA)(1730), Stirling's in his \emph{Methodus Differentialis} (1730).

  [...] Three years later, de Moivre (1733) simplified his result for $p=\tfrac 12$ and showed that the normal density function may be used as an approcimation to the binomial. The generalization to an arbitrary value of $p$ is of course very easy, so without proof de Moivre stated that

  \begin{align*}
      b(np+d,n,p) \sim (2\pi npq)^{-\tfrac 12} \exp\left( -\frac{d^2}{2npq} \right), d = O(\sqrt n),
  \end{align*}

  and that $P_d$ [with $P_d=\P{|\Bs -np|\le d}$] may be obtained by integration. He also showd how to calculate the standardized normal probability integral [numerically] and gave the result for one, two, and three times the standard deviation.
\end{quotation}

\section{Proof of the De Moivre-Laplace theorem}

\section{Applications of De Moivre-Laplace theorem}

Graunt + sex ratio, Computerprogramme

  \chapter{Standard proof of the De Moivre-Laplace theorem}

In this chapter I want to prove the theorem by de Moivre and Laplace whereby I want to follow the basic ideas of their original proof\footnote{You can find similar proofs without error estimations in \cite[pp. 64-67]{irle}, \cite[pp. 131-134]{georgii} and \cite{wiki:demoivrelaplace}}. So we will first derive Stirling's formula of the factorial with which we will find an approximation of the binomial's probability mass function. This approximation will be used to deduce de Moivre's and Laplace's theorem.

\section{Stirling's formula}

Without proof the following theorem\footnote{cf.~\cite[pp. 505, 63]{heuser}} will be used in this section:

\begin{theorem}[Wallis' product]
  It is:

  \begin{align}
    \lim_{n\to\infty} \frac{2^2\cdot4^2\cdot6^2\dots(2n)^2}{1^2\cdot3^2\cdot5^2\dots(2n-1)^2}\cdot \frac{1}{2n} = \lim_{n\to\infty} \frac{1}{2n} \cdot \frac{2^{4n}}{\binom{2n}{n}^2} = \frac{\pi}{2}
  \end{align}
\end{theorem}

\noindent A proof of this equality can be found in~\cite[pp. 504-505]{heuser}. From this equation directly follows:

\begin{align} \label{wallis}
  \lim_{n\to\infty} \frac{2^{2n}}{\sqrt n \binom{2n}{n}} = \lim_{n\to\infty} \sqrt{\frac 1n \cdot \frac{2^{4n}}{\binom{2n}{n}^2}} = \sqrt \pi
\end{align}

\noindent Besides we will also need the Euler–Maclaurin formula~\cite[p. 226]{koenigsberger}:

\begin{theorem}[Euler–Maclaurin formula]
  Let $f:[1,n]\to\R$ be a $2n+1$-times continuously differentiable function. Let $\bn$ be the $n$th Bernoulli number and let $\bp{x}$ be the $n$th periodic Bernoulli polynomial. It is

  \begin{multline} \label{euler-maclaurin-formula}
    \sum_{k=1}^n f(k) = \int_1^n f(x)\d{x} + \frac{f(1)+f(n)}{2} + \sum_{k=1}^n \frac{\bn[2k]}{(2k)!} \Big[f^{(2k-1)}(n) - f^{(2k-1)}(1)\Big] \nl
     + \int_1^n \frac{\bp[2n+1]{x}}{(2n+1)!} f^{(2n+1)}(x) \d{x}
  \end{multline}
\end{theorem}

The periodic Bernoulli polynomials $\bps$ are defined on $[0,1)$ by the following properties~\cite[p. 291]{koenigsberger}:

\begin{enumerate}
  \item $\bp[0]{x} = 1$
  \item $\bps[n+1]'(x) = (n+1)\cdot \bp{x}$
  \item $\int_0^1 \bp{x} \d{x} = 0$
\end{enumerate}

The definition of $\bps$ on $[0,1)$ is then periodically continued to the whole domain $\R$~\cite[pp.~225, 291]{koenigsberger}. The Bernoulli numbers $\bn$ fulfill $\bn=\bp{0}$~\cite[p.~290]{koenigsberger}. All odd Bernoulli numbers $\bn[2n-1]$ for $n>1$ are zero and for the even Bernoulli numbers we get~\cite[p.~289]{koenigsberger}:

\begin{align}
  \bn[2] = \frac 16;\quad \bn[4]=-\frac{1}{30};\quad \bn[6]=\frac{1}{42};\quad \ldots
\end{align}

You can find a proof of the Euler-Maclaurin formula in~\cite[pp.~225-226]{koenigsberger} and~\cite[pp.~506-509]{heuser}. With the above two theorems we can derive Stirling's formula~(cf.~\cite[p.~228]{koenigsberger}):

\begin{theorem}[Stirling's formula]
  The factorial $n!$ fulfills:

  \begin{align}
    n! \in \sqrt{2\pi n} \left(\frac ne\right)^n \fexp{\frac 1{12n} + \bigabs{\frac{1}{120n^2}} }
  \end{align}

  \noindent Therefore

  \begin{align}
     n! \in \sqrt{2\pi n} \left(\frac ne\right)^n \fexp{\bigo{\frac 1n}}
  \end{align}
\end{theorem}

\begin{proof}
  In the proof we will follow~\cite[pp. 227-228]{koenigsberger}. First we apply the Euler-Maclaurin formula~\eqref{euler-maclaurin-formula} to $\ln(n!)=\sum_{k=1}^n \ln(k)$:

  \begin{align}
    \ln(n!) & = \sum_{k=1}^n \ln(k) \nl
    & = \int_1^n \ln(x) \d{x} + \frac{\ln(1)+\ln(n)}{2} + \frac{\bn[2]}{2!} \Big[\ln'(n)-\ln'(1)\Big] + \int_1^n \frac{\bp[3]{x}}{3!} \ln^{(3)}(x) \d{x} \nl
    &= \int_1^n \ln(x) \d{x} + \frac 12 \ln(n) + \frac{1}{12} \left(\frac 1n - 1\right) + \frac 13 \int_1^n \frac{\bp[3]{x}}{x^3} \d{x} \nl
    & \begin{comment}
      \int_1^n \ln(x) \d{x} = \Big[x\ln(x)-x\Big]_1^n = n\ln(n)-n+1
    \end{comment} \nonumber \nl
    &= n\ln(n)-n+\frac 12 \ln(n) + \frac{1}{12n} + \frac{11}{12} + \frac{1}{3} \int_1^n \frac{\bp[3]{x}}{x^3}\d{x} 
  \end{align}

  \noindent Thus

  \begin{align}
    \ln(n!) - n\ln(n) + n - \frac 12 \ln(n) = \frac{11}{12} + \frac{1}{12n} + \frac 13 \int_1^n \frac{\bp[3]{x}}{x^3} \d{x}
  \end{align}

  $\bps[3]$ is on $[0,1)$ as part of a polynomial bounded. Because $\bps[3]$ is $1$-periodic, $\bps[3]$ is bounded on the whole domain $\R$. Thus $\int_1^\infty \frac{\bp[3]{x}}{x^3} \d{x}$ exists because $\int_1^\infty \frac 1{x^3} d{x} < \infty$. Now we define $\seq{b_n}$ via

  \begin{align}
    b_n = \frac{n!}{n^n e^{-n} \sqrt{n}}
  \end{align}
  This sequence converges because
  \begin{align}
    \lim_{n\to\infty} \ln(b_n) & = \lim_{n\to\infty} \Big( \ln(n!) -n \ln(n) + n - \frac 12\ln(n) \Big) \nl
    & = \lim_{n\to\infty} \left( \frac{11}{12} + \frac{1}{12n} + \frac 13\int_1^n \frac{\bp[3]{x}}{x^3} \d{x}\right) \nl
    & = \frac{11}{12} + \frac 13\int_1^{\infty} \frac{\bp[3]{x}}{x^3} \d{x}
  \end{align}
  Let $b=\lim_{n\to\infty} b_n$. To calculate this limit we investigate $\tfrac{b_n^2}{b_{2n}}$:
  \begin{align}
    \lim_{n\to\infty} \frac{b_n^2}{b_{2n}} &= \lim_{n\to\infty} \frac{(n!)^2}{n^{2n}e^{-2n} n} \cdot \frac{(2n)^{2n} e^{-2n} \sqrt{2n}}{(2n)!} \nl
    &= \lim_{n\to\infty} \sqrt{2} \frac{2^{2n}}{\sqrt{n}\binom{2n}{n}} \nl
    & \begin{comment}
    \lim_{n\to\infty} \frac{2^{2n}}{\sqrt{n}\binom{2n}{n}} = \sqrt{\pi} \text{, see \eqref{wallis}}
    \end{comment} \nonumber \nl
    & = \sqrt{2\pi}
  \end{align}

  \noindent On the other hand we have

  \begin{align}
    \lim_{n\to\infty} \frac{b_n^2}{b_{2n}} = \frac{b^2}{b} = b
  \end{align}

  \noindent Thus $\lim_{n\to\infty} b_n = \sqrt{2\pi}$. From this we can follow

  \begin{align}
    \lim_{n\to\infty} \frac{n!}{\sqrt{2\pi n}n^n e^{-n}} = \lim_{n\to\infty} \frac{b_n}{\sqrt{2\pi}} = 1
  \end{align}

  This proves that $n!$ can be approximated by $\sqrt{2\pi n}n^n e^{-n}$ for large $n$. To estimate the error of this approximation we first calculate $\ln(n!)-\ln\left(\sqrt{2\pi n}n^n e^{-n}\right)$:

  \begin{align}
    \ln\left(\frac{n!}{\sqrt{2\pi n}n^n e^{-n}}\right) & = \ln\left(\frac{b_n}{\sqrt{2\pi}}\right) \nl
    &= \ln(b_n)-\ln\left(\sqrt{2\pi}\right) \nl
    &= \ln(b_n)-\lim_{n\to\infty} \ln(b_n) \nl
    &\begin{comment}
      \begin{aligned}
        \ln(b_n) &= \ln(n!)-n\ln(n)+n-\frac 12\ln(n) \nl
        &= \frac{11}{12} + \frac{1}{12n} + \frac 13 \int_1^n \frac{\bp[3]{x}}{x^3} \d{x}
      \end{aligned}
    \end{comment} \nl
    &= \frac{11}{12} + \frac{1}{12n} + \frac 13 \int_1^n \frac{\bp[3]{x}}{x^3} \d{x} -\left(\frac{11}{12}  + \frac 13 \int_1^\infty \frac{\bp[3]{x}}{x^3} \d{x}\right) \nl
    &= \frac{1}{12n}-\frac 13\int_n^\infty \frac{\bp[3]{x}}{x^3} \d{x}
  \end{align}

  On $[0,1)$ the periodic Bernoulli polynomial has the form $\bp[3]{x} = x^3-\frac 32 x^2 + \frac 12 x$~\cite[p.~290]{koenigsberger}. It has roots at $x=0$ and $x=1$. Its derivative $\bps[3]'(x) = 3x^2-3x+\tfrac 12$ has two distinct roots at $x = \frac 12 \pm \sqrt{\frac 1{12}}$. Thus $\bps[3]$ has local extrema at $x=\frac 12 \pm \sqrt{1{12}}$. Because $\bp[3]{0}=\bp[3]{1}=0$ and $\bps[3]$ is 1-periodic its supremum is

  \begin{align}
    \snorm{\bps[3]} = \max\left\{\abs{\bp[3]{\frac 12 + \sqrt{\frac{1}{12}}}}, \abs{\bp[3]{\frac 12 - \sqrt{\frac{1}{12}}}}\right\} \approx 0.0481 \le 0.05 = \frac{1}{20}
  \end{align}

  \noindent Thus

  \begin{align}
    \abs{\frac 13\int_n^\infty \frac{\bp[3]{x}}{x^3} \d{x}} & \le \frac 13 \int_n^\infty \frac{\abs{\bp[3]{x}}}{x^3} \d{x} \nl
    & \le \frac{\snorm{\bps[3]}}{3} \int_n^\infty \frac 1{x^3} \d{x} \nl
    & \le \frac{1}{60} \left[-\frac{1}{2x^2}\right]_n^\infty \nl
    & = \frac 1{120n^2}
  \end{align}

  \noindent So

  \begin{align}
    & \ln\left(\frac{n!}{\sqrt{2\pi n}n^n e^{-n}}\right) = \frac 1{12n} - \frac 13 \int_n^\infty \frac{\bp[3]{x}}{x^3} \d{x} \in \frac 1{12n} + \bigabs{\frac{1}{120n^2}} \bnl
    \implies & n! \in \sqrt{2\pi n} \left(\frac ne\right)^n \fexp{\frac 1{12n} + \bigabs{\frac 1{120n^2}}} \bnl
    &
    \begin{comment}
      \frac 1{12n} + \bigabs{\frac 1{120n^2}} \subseteq \bigo{\frac 1n}
    \end{comment} \bnl
    \implies & n! \in \sqrt{2\pi n} \left(\frac ne\right)^n \fexp{\bigo{\frac 1n}}
  \end{align}
\end{proof}

\section{Local version of de Moivre-Laplace theorem}

In this section we will show that the probability mass function of the standardized binomial distribution can be approximated by the density function of the normal distribution. This will be the basis for the global version of de Moivre-Laplace theorem. The following proof is based on the ideas given by Iosif Pinelis\cite{iosif}.

\begin{theorem}[Local version of de Moivre-Laplace theorem]
  Let $\pol{\x}$ be the set of all polynomials in $\x$. Let $\seq{c_n}$ be a sequence with $\seq{c_n} \in \littleo{\frac 1h}$. For all $\x$ with $0\le k\le n$ and $\abs{x_k} \le c_n$ we have
  \begin{align}
    \bb{x_k} &= \binom nk p^k q^{n-k} \nl
    &\in \fphi{\x} h +\frac{q-p}6 \left(\x^3-3\x\right) \fphi{\x} h^2 + \bigo{\fphi{\x}\pol{\x}h^3}
  \end{align}
\end{theorem}

\begin{proof}
  We have
  \begin{align}
    k &= np + \x \sqrt{npq} \nl
    &= np \left( 1 + \x \frac{\sqrt{q}}{\sqrt{np}}\right) \nl
    &= np \left( 1 + \x q h\right)
  \end{align}
  For $n-k$ we get
  \begin{align}
    n-k &= n-np-\x\sqrt{npq} \nl
    &= nq -\x \sqrt{npq} \nl
    &= nq \left( 1 - \x \frac{\sqrt p}{\sqrt{nq}}\right) \nl
    &= nq \left( 1 - \x p h\right)
  \end{align}
  Because $\x h \in \littleo{1}$ we have
  \begin{align}
    k & \in \bigo{np} \nl
    n-k & \in \bigo{nq}
  \end{align}
  Because $\seq{1+a_n} \in \fexp{a_n + \bigo{a_n^2}}$ we can also conclude
  \begin{align}
    k &\in np \cdot \fexp{\x q h + \bigo{\pol{\x}h^2}} \nl
    n-k &\in nq \cdot \fexp{\x p h + \bigo{\pol{\x}h^2}}
  \end{align}
  Now we can begin with the approximation
  \begin{align}
    \bb{\x} &= \binom nk p^k q^{n-k} \nl
    &= \frac{n!}{k!(n-k)!} p^k q^{n-k} \bnl
    & \begin{comment} \text{Stirling's formula} \end{comment} \bnl
    &\in \frac{\sqrt{2\pi n}n^n e^{-n} \fe{\bigo{\frac 1n}}}{\sqrt{2\pi k}k^k e^{-k} \fe{\bigo{\frac 1k}}\cdot \sqrt{2\pi (n-k)}(n-k)^{n-k} e^{-(n-k)}\fe{\bigo{\frac 1{n-k}}}} \cdot p^k q^{n-k} \bnl
    & \subseteq \sqrt{\frac{n}{2\pi k (n-k)}} \frac{n^n}{k^k(n-k)^{(n-k)}}\cdot p^k q^{n-k} \fe{\bigo{\frac 1n} + \bigo{\frac 1k} + \bigo{\frac 1{n-k}}} \bnl
    &
    \begin{comment}
      \begin{array}{rll}
        \frac 1k & \in \bigo{\frac 1{np}} & \subseteq \bigo{h^2} \bnl
        \frac 1{n-k} & \in \bigo{\frac 1{nq}} & \subseteq \bigo{h^2} \bnl
        \frac 1n & \in \bigo{h^2}
      \end{array}
    \end{comment} \bnl
    & \subseteq \sqrt{\frac{n}{2\pi k (n-k)}} \left(\frac{np}{k}\right)^k \left(\frac{nq}{n-k}\right)^{n-k} \fe{\bigo{h^2}}
  \end{align}

  Now we first approximate the two main factors of the remaining product. We start with the root $\sqrt{\frac{n}{2\pi k(n-k)}}$:
  \begin{align}
    \sqrt{\frac{n}{2\pi k(n-k)}} & \in \sqrt{\frac{n}{2\pi \cdot np \fexp{\x qh +\bigo{\pol{x}h^2}} \cdot nq \fexp{-\x ph + \bigo{\pol{\x}h^2}}}} \bnl
    & \subseteq \sqrt{\frac{1}{2\pi npq \fexp{(q-p)\x h +\bigo{\pol{x}h^2}}}} \bnl
    & \subseteq \frac{1}{\sqrt{2\pi npq}} \fexp{\frac{p-q}2 \x h + \bigo{\pol{\x}h^2}}
  \end{align}

  \noindent In order to approximate $\left(\frac{np}{k}\right)^k \left(\frac{nq}{n-k}\right)^{n-k}$ we investigate its logarithm

  \begin{align}
    \ln\left(\left(\frac{np}{k}\right)^k \left(\frac{nq}{n-k}\right)^{n-k}\right) & = k \ln\left(\frac{np}{k}\right) + (n-k) \ln\left(\frac{nq}{n-k}\right) \nl
    &=-n \left[ \frac kn \ln\left(\frac{k}{np}\right) + \left(1-\frac kn\right) \ln\left(\frac{1-\frac kn}{q}\right)\right] \nl
    &
    \begin{comment}
      \h{s} := -n\left[s \ln\left(\frac sp\right) + (1-s) \ln\left(\frac{1-s}{1-p}\right)\right]
    \end{comment} \nl
    &= \h{\frac kn}
  \end{align}

  \noindent The function $\h{s}$ has the derivatives

  \begin{align}
    h'(s) &= -n\left[\ln\left(\frac sp\right) - \ln\left(\frac{1-s}{1-p}\right)\right] \nl
    h''(s) &= -\frac ns - \frac n{1-s} \nl
    h^{(3)}(s) &= \frac n{s^2} - \frac n{(1-s)^2} \nl
    h^{(4)}(s) &= -\frac {2n}{s^3} - \frac {2n}{(1-s)^3}
  \end{align}

  \noindent We have

  \begin{align}
    h(p)& = 0 \nl
    h'(p)& =0 \nl
    h''(p) & = -\frac np - \frac nq = -\frac n{pq} \nl
    h^{(3)}(p) &= \frac n{p^2}-\frac{n}{q^2} = \frac{n\left(q^2-p^2\right)}{p^2q^2} \nl
    &
    \begin{comment}
      q^2-p^2 = (q-p)(q+p) = q-p
    \end{comment} \nl
    &= -\frac{n(p-q)}{p^2q^2}
  \end{align}

  \noindent Therefore the Taylor series of $\h{s}$ is

  \begin{align}
    h(s) = -\frac{n}{2pq} (s-p)^2 - \frac{n(p-q)}{6p^2q^2} (s-p)^3 + R(s)
  \end{align}

  with $R(s)=-\frac{2n}{4!} \left(\frac{1}{\tilde s^3}+\frac{1}{(1-\tilde s)^3}\right) (s-p)^4$ for a $\tilde s$ in the interval bounded by $s$ and $p$. We want to evaluate this series for $s=\frac kn$. Here we have $s-p=\x pqh$ and thus

  \begin{align}
    \begin{array}{rll}
      -\frac{n}{2pq} (s-p)^2 &= -\frac{n}{2pq} \x^2 p^2 q^2 h^2 & = -\frac{\x^2}{2} \nl
      -\frac{n(p-q)}{6p^2q^2} (s-p)^3 &= -\frac{n(p-q)}{6p^2q^2} \x^3p^3q^3h^3 & = -\frac{p-q}{6} \x^3 h
    \end{array}
  \end{align}

  Because $\tilde s$ is in the interval bounded by $p$ and $s=\frac kn \in p \cdot \bigo{1} \subseteq \bigo{1}$ we deduce $\tilde s \in \bigo{1}$, too. Similarly we get $1-\tilde s\in \bigo{1}$ because $1-s = \frac{n-k}{n} \in q \cdot \bigo{1}$. From this follows:

  \begin{align}
    \frac 1{\tilde s^3} + \frac 1{(1-\tilde s)^3} & \in \bigo{1}
  \end{align}

  \noindent So

  \begin{align}
    R(s) & = -\frac{2n}{4!} \left(\frac 1{\tilde s^3} + \frac 1{(1-\tilde s)^3}\right) (s-p)^4 \nl
    & \in -\frac{2n}{4!} \bigo{1} \x^4 p^4 q^4 h^4 \nl
    &
    \begin{comment}
      n h^2 \in \bigo{1}
    \end{comment} \nl
    & \subseteq \bigo{\pol{\x} h^2}
  \end{align}

  \noindent Thus

  \begin{align}
    \h{\frac kn} &=-\frac{n}{2pq} (s-p)^2 - \frac{n(p-q)}{6p^2q^2} (s-p)^3 + R(s) \nl
    & \in -\frac{\x^2}{2} - \frac{p-q}{6} \x^3 h + \bigo{\pol{\x} h^2}
  \end{align}

  \noindent Finally

  \begin{align}
    \left(\frac{np}{k}\right)^k \left(\frac{nq}{n-k}\right)^{n-k} & = \fexp{\h{\frac kn}} \bnl
    & \in \fexp{-\frac{\x^2}{2} - \frac{p-q}{6} \x^3 h + \bigo{\pol{\x} h^2}} \bnl
    & \subseteq \fexp{-\frac{\x^2}{2}} \fexp{-\frac{p-q}6 \x^3 h + \bigo{\pol{\x} h^2}}
  \end{align}

  \noindent Now we can complete the approximation

  \begin{align}
    \bb{\x} & \in \sqrt{\frac{n}{2\pi k(n-k)}} \left(\frac {np}{k}\right)^k \left(\frac{nq}{n-k}\right)^{n-k} \fe{\bigo{h^2}} \bnl
    &
    \begin{comment}
      \begin{aligned}
        \sqrt{\frac{n}{2\pi k(n-k)}} & \in \frac1{\sqrt{2\pi npq}} \fexp{\frac{p-q}2 \x h + \bigo{\pol{\x}h^2}}\bnl
        \left(\frac {np}{k}\right)^k \left(\frac{nq}{n-k}\right)^{n-k} & \in \fexp{-\frac{\x^2}2} \fexp{-\frac{p-q}6 \x^3 h + \bigo{\pol{\x}h^2}}
      \end{aligned}
    \end{comment} \bnl
    & \subseteq \frac 1{\sqrt{2\pi npq}} \cdot \fexp{-\frac{\x^2}{2}} \fexp{\frac{p-q}2 \left(\x  - \frac 13 \x^3\right) h + \bigo{\pol{\x}h^2}} \bnl
    &
    \begin{comment}
      \fexp{a_n} \in \left(1+a_n+\bigo{a_n^2}\right)
    \end{comment} \bnl
    & \subseteq \frac 1{\sqrt{2\pi npq}} \cdot \fexp{-\frac{\x^2}{2}} \left(1 +\frac{p-q}2 \left(\x - \frac 13 \x^3 \right) h + \bigo{\pol{\x}h^2}\right) \bnl
    &
    \begin{comment}
      \fphi{x} = \frac{1}{\sqrt{2\pi}}\fexp{-\frac{x^2}2} \land h = \frac{1}{\sqrt{npq}}
    \end{comment} \bnl
    & \subseteq \fphi{\x} h \left(1 +\frac{q-p}6 \left(\x^3-3\x \right) h + \bigo{\pol{\x}h^2}\right) \bnl
    & \subseteq \fphi{\x} h +\frac{q-p}6 \left(\x^3-3\x\right) \fphi{\x} h^2 + \bigo{\fphi{\x}\pol{\x}h^3}
  \end{align}
\end{proof}

\section{First global version of de Moivre-Laplace}

\begin{theorem}[First global version of de Moivre-Laplace]
  Let $a,b\in \R$ with $a \le b$. We have
  \begin{align}
    \P{a \le \BBs \le b} & = \sum_{a \le \x \le b} \bb{\x} \nl
    &
    \begin{multlined}
      \in \int_a^b \fphi{t} \d{t} + \frac{q-p}6 \left[\br{b^2-1}\fphi{b}-\br{a^2-1}\fphi{a}\right] h\nl
      +\bigabs{\snorm{\phi}h}+ \bigo{h^2}
    \end{multlined}
  \end{align}
\end{theorem}

\todo{Vorzeichenfehler}

\begin{remark}
  With the above theorem we can establish the same approximation Uspensky found in his textbook ``Introduction to mathematical probability'' of 1937 where he stated \cite[p. 129]{uspensky}
  \begin{align}
    \P{a+\frac 12h \le \BBs \le b-\frac 12h} = \int_a^b \fphi{x} \d{x} + \frac{q-p}6 \br{\br{1-b^2}\fphi{b} - \br{1-a^2}\fphi{a}} h + \epsilon
  \end{align}
  whereby the error term $\epsilon$ fulfills for $npq \ge 25$ the inequality
  \begin{align} \label{uspensky:error}
    \abs{\epsilon} < \frac{0.13+0.18\abs{p-q}}{npq} + \fexp{-\frac 32 \sqrt{npq}}
  \end{align}
  Thereby Uspensky uses the characteristic function of the binomial distribution to prove his theorem \cite[pp. 119-129]{uspensky}. With the following proof we show that the same approximation is possible with Stirling's formula which demonstrates the power of this proof method. However we cannot establish the same error bound Uspensky found with \eqref{uspensky:error}.
\end{remark}

\begin{proof}
  In the following we will only consider those $n\in\N$ with $\x[0] \le a \le b \le \x[n]$. Because $\lim_{n\to\infty} \x[0] = -\infty$ and $\lim_{n\to\infty} \x[n] = \infty$ these are almost all $n\in\N$. Because $c_n=\max\{\abs{a},\abs{b}\}$ is a sequence with $c_n\in\littleo{\frac 1h}$ and for all $\x$ with $a \le \x \le b$ we have $\abs{\x} \le c_n$ we can apply the local version of de Moivre-Laplace and we get

  \begin{align}
    \sum_{a \le \x \le b} \bb{\x} \in \sum_{a \le \x \le b} \br{\fphi{\x} h +\frac{q-p}6 \left(\x^3 - \x\right) \fphi{\x} h^2 + \bigo{\fphi{\x}\pol{\x}h^3}}
  \end{align}

  We have three Riemann sums. For any $\f{x}=\fphi{x}p(x)$ with a polynomial $p(x)$ we have $\snorm{f} < \infty$ and $\snorm{f'}<\infty$. So we apply the following rule from theorem \ref{thm:arithmetic_rules}

  \begin{align}
    \sum_{a \le \x \le b} \f{\x} h & \in \int_a^b \f{t} \d{t} + \bigabs{\snorm{f}h + \frac 14 \snorm{f'} (b-a) h + \frac 14 \snorm{f'} h^2} \bnl
    & \subseteq  \int_a^b \f{t} \d{t} + \bigo{h}
  \end{align}
  Thus we have for the third sum
  \begin{align}
    \sum_{a \le \x \le b} \bigo{\fphi{\x} \pol{\x} h^3} & \subseteq \bigo{\int_a^b \fphi{t} \pol{t} \d{t} \cdot h^2 + \bigo{h^3}} \subseteq \bigo{h^2}
  \end{align}
  For the second term we get
  \begin{align}
    \sum_{a \le \x \le b} \frac{q-p}6 \br{\x^3 - x} \fphi{\x} h^2 &\in \frac{q-p}6 \int_a^b \br{t^3-t}\fphi{t} \d{t} \cdot h + \bigo{h^2} \bnl
    &= \frac{q-p}6 \left[ \br{t^2-1}\fphi{t} \right]_a^b h + \bigo{h^2} \bnl
    &= \frac{q-p}6 \left[ \br{b^2-1}\fphi{b} - \br{a^2-1}\fphi{a} \right] h + \bigo{h^2}
  \end{align}
  For $\sum_{a\le b\le b} \fphi{\x}h$ we need a better approximation. From the following equation
  \begin{align}
    \int_{\x-\frac h2}^{\x+\frac h2} \fphi{t} \d{t} & = \int_{-\frac h2}^{+\frac h2} \frac{\fphi{\x+t}+\fphi{\x-t}}2 \d{t} \bnl
    & = \fphi{\x} \int_{-\frac h2}^{\frac h2} \fcosh{\x t} \fexp{-\frac{t^2}2} \d{t} \bnl
    &
    \begin{comment}
      \fcosh{\x t}\fexp{-\frac{t^2}2} = 1 + \bigo{t^3}
    \end{comment} \bnl
    & \in \fphi{\x} h + \bigo{h^3}
  \end{align}

  \noindent we can follow

  \begin{align}
    \fphi{\x} h \in \int_{\x-\frac h2}^{\x+\frac h2} \fphi{t} \d{t} + \bigo{h^3}
  \end{align}

  \noindent and thus (let $\tilde a=\xup{a}$ and $\tilde b=\xdown{b}$)

  \begin{align}
    \sum_{a \le \x \le b} \fphi{\x} h &\in \sum_{a \le \x \le b} \br{\int_{\x-\frac h2}^{\x+\frac h2} \fphi{t} \d{t} + \bigo{h^3}} \bnl
    &\subseteq \int_{\tilde a - \frac h2}^{\tilde b+\frac h2} \fphi{t} \d{t} + \bigo{h^2} \bnl
    &\subseteq \int_a^b \fphi{t} \d{t} + \bigabs{\snorm{\phi} h} + \bigo{h^2}
  \end{align}

  \noindent So we get

  \begin{align}
    \sum_{a\le \x \le b} \bb{\x} &\in \sum_{a \le \x \le b} \left(\fphi{\x} h +\frac{q-p}6 \left(\x^3 - \x\right) \fphi{\x} h^2 + \bigo{\fphi{\x}\pol{\x}h^3}\right) \bnl
    &
    \begin{comment}
      \begin{aligned}
        \sum_{a \le \x \le b} \fphi{\x} h & \in \int_a^b \fphi{t} \d{t} + \bigabs{\snorm{\phi} h} + \bigo{h^2} \bnl
        \sum_{a \le \x \le b} \frac{q-p}6 \br{\x^3-\x} \fphi{\x} h^2 & \in \frac{q-p}6 \int_a^b \br{t^3-t} \fphi{t} \d{t}\cdot h + \bigo{h^2} \bnl
        \sum_{a \le \x \le b} \bigo{\fphi{\x}\pol{x}h^3} &\subseteq \bigo{h^2}
      \end{aligned}
    \end{comment} \bnl
    &
    \begin{multlined}
      \subseteq \int_a^b \fphi{t} \d{t} + \frac{q-p}6 \left[\br{b^2-1}\fphi{b}-\br{a^2-1}\fphi{a}\right] h \nl
      + \bigabs{\snorm{\phi}h} + \bigo{h^2}
    \end{multlined}
  \end{align}
\end{proof}

\section{Second global version of the de Moivre-Laplace theorem}

\begin{theorem}[Second global version of the de Moivre-Laplace theorem]
  For any $z\in\R$ we find

  \begin{align}
    \P{\BBs \le z} &= \sum_{\x \le z} \bb{\x} \bnl
    & \in \int_{-\infty}^z \fphi{t} \d{t} + \frac{q-p}6 \br{z^2-1}\fphi{z} h + \littleo{h}
  \end{align}
\end{theorem}

\begin{proof}
  We consider only those $n\in\N$ with $\x[0] \le z \le \x[n]$. Because $\lim_{n\to\infty} \x[0] = -\infty$ and $\lim_{n\to\infty} \x[n] = \infty$ these are almost all $n\in\N$. Let $\seq{c_n}$ be a sequence with
  \begin{itemize}
    \item $\x[0] \le c_n \le \min\{z,0\}$
    \item $\seq{\abs{c_n}}\in\littleo{\frac 1h}$
    \item $\seq{\frac{1}{c_n^2}} \in \littleo{h}$
    \item $\lim_{n\to\infty} c_n=-\infty$
  \end{itemize}
  For example one may take $c_n = \min\left\{z, 0, \xdown{- \sqrt[3]{n}\sqrt{\frac pq}}\right\}$. We split the sum $\sum_{\x \le z} \bb{\x}$ into two sums:

  \begin{align}
    \sum_{\x \le z} \bb{\x} = \sum_{\x < c_n} \bb{\x} + \sum_{c_n \le \x \le z} \bb{\x}
  \end{align}
  The first sum can be estimated with the Chebyshev's inequality:
  \begin{align}
    0 \le \sum_{\x < c_n} \bb{\x} \le \P{\abs{\BBs} > \abs{c_n}} \le \frac{1}{c_n^2}
  \end{align}
  Similarly we get
  \begin{align}
    0 \le \int_{t < c_n} \fphi{t}\d{t} \le \int_{\abs{t} > \abs{c_n}} \fphi{t} \d{t} \le \frac{1}{c_n^2}
  \end{align}
  Thus
  \begin{align}
    \sum_{\x < c_n} \bb{\x} & \in \int_{t < c_n} \fphi{t}\d{t} + \bigabs{\frac{1}{c_n^2}} \nl
    & \subseteq \int_{t < c_n} \fphi{t}\d{t} + \littleo{h}
  \end{align}

  The second sum can be approximated as in the above theorem. Because $\abs{\x} \le \max\{\abs{c_n},z\} \in \littleo{\frac 1h}$ we can apply the local version of de Moivre-Laplace and we get
  \begin{align}
    \sum_{c_n \le \x \le z} \bb{\x} = \sum_{c_n \le \x \le z} \br{\fphi{\x} h +\frac{q-p}6 \left(\x^3 - \x\right) \fphi{\x} h^2 + \bigo{\fphi{\x}\pol{\x}h^3}}
  \end{align}
  Now we have for any $f(x)=\fphi{x}p(x)$ with $p(x)$ being a polynomial:
  \begin{align}
    \sum_{c_n \le \x \le b} \f{\x} h & \in \int_{c_n}^b \f{t} \d{t} + \bigabs{\snorm{f}h + \frac 14 \snorm{f'} (b-c_n) h + \frac 14 \snorm{f'} h^2} \nl
    &
    \begin{comment}
      (b-c_n)h \in \littleo{1}
    \end{comment} \nl
    & \subseteq  \int_{c_n}^b \f{t} \d{t} + \littleo{1}
  \end{align}
  Therefore
  \begin{align}
    \sum_{c_n \le \x \le z} \bigo{\fphi{\x}\pol{\x}h^3} \in \bigo{\int_{c_n}^z \fphi{t} \d{t} \cdot h^2 + \littleo{h^2}} \subseteq \littleo{h}
  \end{align}
  and
  \begin{align}
    \sum_{c_n \le \x \le z} \frac{q-p}6 \br{\x^3-\x} \fphi{\x} h^2 &\in \frac{q-p}6 \int_{c_n}^z \br{t^3-t} \fphi{t} \d{t} \cdot h + \littleo{h} \bnl
    &
    \begin{multlined}
      = \frac{q-p}6 \int_{-\infty}^z \br{t^3-t} \fphi{t} \d{t} \cdot h \bnl
      - \frac{q-p}6 \int_{-\infty}^{c_n} \br{t^3-t} \fphi{t} \d{t} \cdot h + \littleo{h}
    \end{multlined} \bnl
    &
    \begin{comment}
      \lim_{n\to\infty} c_n = -\infty \implies \lim_{n\to\infty} \int_{-\infty}^{c_n} \br{t^3-t} \fphi{t} \d{t} = 0
    \end{comment} \bnl
    & = \frac{q-p}6 \int_{-\infty}^z \br{t^3-t} \fphi{t} \d{t} \cdot h + \littleo{h} \bnl
    & = \frac{q-p}6 \left[\br{t^2-1} \fphi{t}\right]_{-\infty}^{z} \cdot h + \littleo{h} \bnl
    & = \frac{q-p}6 \br{z^2-1} \fphi{z} \cdot h + \littleo{h}
  \end{align}
  As in the previous proof we have
  \begin{align}
    \sum_{c_n \le \x \le z} \fphi{\x} h &\in \sum_{c_n \le \x \le z} \br{\int_{\x-\frac h2}^{\x+\frac h2} \fphi{t}\d{t} + \bigo{h^3}} \bnl
    &
    \begin{comment}
      (z-c_n)h \in \littleo{1}
    \end{comment} \bnl
    &\subseteq \int_{c_n}^z \fphi{t} \d{t} + \bigabs{\snorm{\phi} h} + \littleo{h} \bnl
    &\subseteq \int_{c_n}^z \fphi{t} \d{t} + \littleo{h}
  \end{align}
  So we can summarize
  \begin{align}
    \sum_{c_n \le \x \le z} \bb{\x} & \in \sum_{c_n \le \x \le z} \br{\fphi{\x} h +\frac{q-p}6 \br{\x^3 - \x} \fphi{\x} h^2 + \bigo{\fphi{\x}\pol{\x}h^3}} \bnl
    &
    \begin{comment}
      \begin{aligned}
        \sum_{c_n \le \x \le z} \fphi{\x} h &\in \int_{c_n}^z \fphi{t} \d{t} + \littleo{h} \bnl
        \sum_{c_n \le \x \le z} \frac{q-p}6 \br{\x^3-\x} \fphi{\x} h^2 &\in \frac{q-p}6 \br{z^2-1} \fphi{z} \cdot h + \littleo{h} \bnl
        \sum_{c_n \le \x \le z} \bigo{\fphi{\x}\pol{\x}h^3} & \in \littleo{h}
      \end{aligned}
    \end{comment} \bnl
    & \subseteq \int_{c_n}^z \fphi{t} \d{t} + \frac{q-p}6 \br{z^2-1}\fphi{z} h + \littleo{h}
  \end{align}

  \noindent So we can finish the proof

  \begin{align}
    \sum_{\x \le z} \bb{\x} & = \sum_{\x < c_n} \bb{\x} + \sum_{c_n \le \x \le z} \bb{\x} \bnl
    & \in \int_{t < c_n} \fphi{t} \d{t} + \int_{c_n}^z \fphi{t} \d{t} + \frac{q-p}6 \br{z^2-1}\fphi{z} h + \littleo{h} \bnl
    & \in \int_{-\infty}^z \fphi{t} \d{t} + \frac{q-p}6 \br{z^2-1}\fphi{z} h + \littleo{h}
  \end{align}
\end{proof}

\section{Conclusion}

\includeplot{edgeworth_diff_pdf}{This plot shows $f(\x) = \bb{\x}-\fphi{\x}h$ and $g(\x)=\frac{q-p}6 \left(\x^3-3\x\right) \fphi{\x} h^2$ for $n=1000$ and $p=\frac 14$. It demonstrates that our correction term $\frac{q-p}6 \left(\x^3-3\x\right) \fphi{\x} h^2$ describes very well the error $\bb{\x}-\fphi{\x}$ in the normal approximation.}

We see that the standard proof of the de Moivre-Laplace theorem with Stirling's formula yield a good estimation of the approximation's error. We do not only get the right order $\bigo{h}$ of the error's convergence as it is shown by the Berry-Esseen theorem, we also derived a better approximation for the binomial distribution which is

\begin{align}
  \P{\BBs = \x}& \approx \fphi{\x} h +\frac{q-p}6 \left(\x^3-3\x\right) \fphi{\x} h^2 \nl
  \P{a \le \BBs \le b} & \approx  \int_a^b \fphi{t} \d{t} + \frac{q-p}6 \left[\br{b^2-1}\fphi{b}-\br{a^2-1}\fphi{a}\right] h \nl
  \P{\BBs \le z} & \approx \int_{-\infty}^z \fphi{t} \d{t} + \frac{q-p}6 \br{z^2-1}\fphi{z} h
\end{align}

These extensions of the normal approximation are precisely the first two terms of the Edgeworth Expansion \cite[pp. 39 ff.]{hall}\cite{wiki:edgeworth}. The Edgeworth Expansion is a series which approximates the cumulative distribution function of a random variable \cite{wiki:edgeworth}. In the case of the binomial distribution the Edgeworth Expansion is
\begin{align}
  \P{\BBs \le z} & \approx \int_{-\infty}^z \fphi{t} \d{t} + \frac{q-p}6 \br{z^2-1}\fphi{z} h + R\left(np + z \sqrt{npq}\right)\fphi{z}h + \bigo{h^2}
\end{align}
with $R(x)=\lfloor x \rfloor -x+\frac 12$ \cite[p. 46]{hall}. By comparing the above equation with our second global version of de Moivre-Laplace we see that we can provide the same first term of the Edgeworth Expansion, but with a worse estimation of the remaining error ($\littleo{h}$ instead of $\bigo{h^2}$). This demonstrates the power of the standard proof with Stirling's formula.

\todo{Folgendes sollte untersucht werden: Vergleich mit numerischen Diagrammen.}


  \chapter{Alternative proof via slope approximation}

In this chapter, I present an alternative proof of the de Moivre-Laplace theorem which approximates the slope of the probability mass function by an ordinary differential equation\footnote{Instead of the term ``ordinary differential equation'' also its acronym ``ODE'' is used~\cite[p. 2]{ricardo}\cite{wiki:ode}.}. The idea of this proof comes from Prof. Peter Pickl and was first described in the thesis ``Der Zentrale Grenzwertsatz'' by Karolina Fels~\cite{fels}. However, the idea to derive an ODE in order to find an approximation for the binomial distribution occurred already in two papers by Daniel Bernoulli in 1770 and 1771~\cite{hald1}.

\section{Intuition behind the proof}

In this section I want to present the idea of the proof of Karolina Fels \cite[pp. 16-19]{fels}. Thereby, I will slightly change her proof sketch. She started with the function $f(x) = \binom nk p^kq^{n-k}$, whereby $k=np+x$, and approximated $f'(x)$ by calculating $f(x)-f(x-1)$~\cite[p. 16]{fels}. We will approximate the slope of $\bbs$ between $\bb{\x[k+1]}$ and $\bb{\x[k]}$:

\begin{align}
  \frac{\bb{\x[k+1]}-\bb{\x[k]}}h &= \frac 1h \br{\binom n{k+1}p^{k+1}q^{n-k-1} - \binom nk p^k q^{n-k}} \nl
  &= \frac 1h \binom nk p^k q^{n-k} \br{\frac{(n-k)p}{(k+1)q}-1} \nl
  &= \frac 1h \bb{\x[k]} \frac{np-kp-kq-q}{(k+1)q} \nl
  &= - \frac 1h \bb{\x[k]} \frac{k-np+q}{(k+1)q} \nl
  &= - \frac 1h \bb{\x[k]} \br{\frac{k-np}{npq} + \frac{q}{npq}} \frac{np}{k+1} \nl
  &= - \frac 1h \bb{\x[k]} \br{\x[k]h + q h^2} \frac{np}{k+1} \nl
  &= -\x[k] \bb{\x[k]} \frac{np}{k+1} + \bb{\x[k]} q h \frac{np}{k+1}
\end{align}

Let us assume that the cumulative distribution of $\BBs$ can be approximated by a distribution with density function $\phi(x)$ for $n\to\infty$. Under the limit $n\to\infty$ we have $h\to 0$. For a fixed $x\in\R$ we also set $\x = \xrnd{x} = \frac{\rnd{np+ x \sqrt{npq}}-np}{\sqrt{npq}}$, which is the closest $\x$ to $x$. We get
\begin{align}
  \frac{\bb{\x[k+1]}-\bb{\x[k]}}h & \to \phi'(x) \nl
  \bb{\x[k]} &\to \fphi{x} \nl
  \frac{np}{k+1} &= \frac{np}{\rnd{np+x\sqrt{np}}+1} \to 1
\end{align}
Therefore $\fphi{x}$ is described by the ODE
\begin{align}
  \phi'(x) = -x\phi(x)
\end{align}

This ODE is solved by $\phi_c(x) = c \fexp{-\frac{x^2}2}$. Because $\int_\R \fphi{1} \d{t} = 1$ and $\int_\R \fexp{-\frac{t^2}2} d{t} = \sqrt{2\pi}$~\cite[p. 47]{georgii} we get finally $\fphi{x} = \frac{1}{\sqrt{2\pi}} \fexp{-\frac{x^2}2}$.

\section{Necessary prior knowledge from ODE theory}

In order to formalize the above proof sketch into a valid proof we will need some concepts and theorems from the theory of ordinary differential equations. An equation of the form

\begin{align} \label{ode:def}
  y'(x) = f(x,y(x))
\end{align}

with $f:I \times \R \to \R$ and $I$ being a connected subset of $\R$ will be called an \emph{ordinary differential equation} (of first order)~\cite[p. 465]{stoer}\cite{wiki:ode}. Any function $y:\R\to\R$ fulfilling the ODE \eqref{ode:def} is called a \emph{solution of the ordinary differential equation}~\cite[p. 8]{ricardo}\cite{wiki:ode}. For example, the density function $\fphi{x} = \frac{1}{\sqrt{2\pi}} \fexp{-\frac{x^2}2}$ of the normal distribution is the solution of the ODE
\begin{align}
  y'(x) = -xy(x)
\end{align}
with the initial value $y(0) = \frac{1}{\sqrt{2\pi}}$. This follows from the theorem:

\begin{theorem} \label{chr_ode}
  The solution to the ODE $y'(x) = -xy(x)$ with the initial value $y(a)=b$ for $a\in\R$ and $b\in \R\setminus\{0\}$ is
  \begin{align}
    y(x) = b \fexp{-\frac{x^2}2+\frac{a^2}2} = \frac{b}{\fphi{a}} \fphi{x}
  \end{align}
\end{theorem}

\begin{proof}
  The ODE $y'(x)=-xy(x)$ can be solved by separation of variables (see \cite[p. 28]{ricardo} and \cite{wiki:separation}). We can derive:

  \begin{align}
      && y'(x)& =-xy(x) \bnl
      &\iff & \int_{y(a)}^{y(x)} \frac 1y \d{y} & = \int_a^x -t \d{t} \bnl
      &\iff & \left[\ln{y}\right]^{y(x)}_{y(a)=b} & = \left[-\frac{t^2}2\right]_a^x \bnl
      &\iff & \ln(y(x))-\ln(b) &= -\frac{x^2}2 + \frac{a^2}2 \bnl
      &\iff & y(x) &= b \fexp{-\frac{x^2}2 + \frac{a^2}2} = \frac b{\fphi{a}}\fphi{x}
  \end{align}
\end{proof}

ODEs can be approximated numerically with \emph{one-step methods}~\cite[pp.~471~ff]{stoer}. A one-step method is described by a function $M(x,y,h,f)$ with a step size $h>0$~\cite[p.~473]{stoer}. Given an initial value $y(x_0)=y_0$ the approximations $y_k$ at the points $x_k = x_0 + kh$ ($k\in\N$) can be obtained recursively via the relation \cite[p. 473]{stoer}

\begin{align}
  y_{k+1} = y_k + M(x_k, y_k, h, f) \cdot h
\end{align}

One of the most famous one-step methods is the \emph{Euler method} with $M(x,y,h,f) = f(x,y)$~\cite[p. 473]{stoer}. We will use the concept of one-step methods to connect the density function of the normal distribution with the probability mass function of the binomial distribution. To state the theorem of the global error of one-step methods we will need the following definition of the local discretization error (see \cite[pp. 473-474]{stoer}):

\begin{definition}[Local discretization error]
  The \emph{local discretization error} is the error of the one-step method which is produced during a step of approximation. Given a one-step method $M(x,y,h,f)$, its local discretization error at the point $x_k$ is the value

  \begin{align}
    \tau(x_k) = \frac{y(x_{k+1})-y(x_k)}h - M(x_k,y(x_k),h,f)
  \end{align}
\end{definition}

\noindent Therefore, $\tau$ fulfills:

\begin{align}
  \underbrace{y(x_{k+1})}_{\text{exact solution}} = \underbrace{y(x_k) + M(x_k,y(x_k),h,f)h}_{\text{approximation of one-step method}} + \underbrace{\tau(x_k)h}_{\text{error}}
\end{align}

\noindent The error of an one-step method can be estimated as in~\cite[pp. 478-479]{stoer}

\begin{theorem}[Global error of one-step methods] \label{thm:ode_global_error}
  Let $f:I\times \R\to\R$ and let $y(x)$ be the solution to the initial value problem

  \begin{align}
    y'(x)=f(x,y(x)),\ y'(x_0)=y_0,\ x_0 \in I
  \end{align}

  Let $M(x,y,h,f)$ be a continuous one-step method for this ODE with step size $h > 0$. Let $M(x,y,h,f)$ be Lipschitz continuous\footnote{See \cite[p. 467]{stoer} and \cite{wiki:lipschitz} for more details about Lipschitz continuity.} regarding the variable $y$ with the Lipschitz constant $L > 0$, i.e. for all $x\in I$ and $y_1,y_2\in\R$ we find
  \begin{align}
    \abs{M(x,y_1,h,f)-M(x,y_2,h,f)} \le L \abs{y_1-y_2}
  \end{align}
  Let the absolute value of the local discretization error be bounded by $C h^p$ for a $C > 0$ and $p\in \N$, i.e. for all $x_k \in I$:
  \begin{align}
    \abs{\tau(x_k)} = \abs{\frac{y(x_{k+1})-y(x_k)}h - M(x_k,y(x_k),h,f)} \le Ch^p
  \end{align}
  Then the global error $\abs{y(x_k)-y_k}$ is bounded by $\frac 1L \br{\fe{L\abs{x_k-x_0}}-1} C h^p$:
  \begin{align}
    \abs{y(x_k)-y_k} \le \frac{\fe{L\abs{x_k-x_0}}-1}{L} \cdot Ch^p
  \end{align}
\end{theorem}

We may also have a one-step backwards method by setting $h < 0$. The above definitions and the  above theorem \ref{thm:ode_global_error} are applicable in this case, too.

\section{Local version of the de Moivre-Laplace theorem}

We will prove the local version of the de Moivre-Laplace theorem in the following form:

\newcommand*{\bnp}{\bb{\xnp}}
\newcommand*{\knp}{m}
\newcommand*{\xnp}{\x[\knp]}
\newcommand*{\phin}[2][n]{y_{#1}\left({#2}\right)}
\newcommand*{\dphin}[2][n]{y_{#1}'\left({#2}\right)}

\begin{theorem}[Local version of the de Moivre Laplace theorem]
  Let $I$ be a bounded interval and let $\xnp$ be the mode of $\bbs$ (i.e. the $\x$ where $\bbs$ attains its maximum).  Then there are $C > 0$ and $L > 0$ independent of $\x$ and $n$ such that

  \begin{align}
    \bb{\x} \in \frac{\bnp}{\fphi{\xnp}}\fphi{\x} + \bigabs{{\bnp}\cdot \frac{\fe{L\abs{\x-\xnp}}-1}L Ch}
  \end{align}
\end{theorem}

\begin{remark} \label{remark:alt_local}
  The mode of $\xnp$ is the $\x[k]$ with $k = \lfloor (n+1)p \rfloor$ (cf. \cite{wiki:binomial_distribution}\cite{nicolas}). Thus we have
  \begin{align}
    np+p-1 \le \knp \le np+p \implies \abs{\knp-np} \le 1 \implies \abs{\xnp} \le h \implies \lim_{n\to\infty} \xnp = 0
  \end{align}
  We may also generalize the theorem to each bounded set $I \subset \R$. Besides because of
  \begin{align}
    \abs{\x[k]-\xnp} \le \abs{\x[k]}+\abs{\xnp} \le L + h \le L + \frac{1}{\sqrt{pq}}
  \end{align}
  we can directly deduce that there is a $\tilde C > 0$ independent of $\x$ and $n$ with
  \begin{align}
    \bb{\x} \in \frac{\bnp}{\fphi{\xnp}}\fphi{\x} + \bigabs{{\bnp}\cdot \tilde Ch}
  \end{align}

  Note that we get the right convergence order $\bigo{h^2}$ of the remaining error because $\bnp \in \bigo{h}$ as we will prove later. However, the error estimation of this theorem is less strong than that of the standard proof with Stirling's formula.
\end{remark}

This proof bases on the theory of one-step methods for approximating ordinary differential equations. Thereby, we will revert the perspective: Whereas in the theory of ODEs one-step methods are approximations of unknown solutions of ODEs, here we will consider the probability mass function of the binomial distribution as the result of an one-step method. This one-step method corresponds to an ODE and the solution of this ODE will be an approximation to the probability mass function. Thus we flip the role of the exact solution and the approximation.

\begin{proof}
  \newcommand*{\Lt}{\tilde L}

  Let $\Lt$ be the maximum of $h$ and $\sup\{\abs{x}: x \in I\}$ and let $\pol{\x}$ be the set of all polynomials in $\x$. We extend $I$ to $\interval{-\Lt}{\Lt}$ so that $\xnp \in I$ because $\Lt \ge h$. We start with investigating $\bb{\x[k+1]}-\bb{\x[k]}$

  \begin{align}
    \bb{\x[k+1]} - \bb{\x[k]} &= \binom n{k+1} p^{k+1}q^{n-k-1} - \binom nk p^k q^k \nl
    &= \binom nk p^k q^{n-k} \br{\frac{\binom n{k+1}p}{\binom nk q} -1} \nl
    &= \bb{\x[k]} \br{\frac{k! (n-k)! p}{(k+1)!(n-k-1)! q} -1} \nl
    &= \bb{\x[k]} \br{\frac{(n-k)p}{(k+1)q}-1} \nl
    &
    \begin{comment}
      k = np + \x[k]\sqrt{npq}
    \end{comment} \nl
    &= \bb{\x[k]} \br{\frac{npq-\x[k]p\sqrt{npq}}{npq + \x[k]q\sqrt{npq}+q}-1} \nl
    &= \bb{\x[k]} \br{\frac{1-\x[k]ph}{1+\x[k]qh+qh^2}-1} \nl
    &
    \begin{comment}
      \frac1{1+x} = 1 - x+ \bigo{x^2}
    \end{comment} \nl
    &\in \bb{\x[k]} \br{\br{1-\x[k]ph}\br{1-\x[k]qh + \bigo{\pol{x}h^2}}-1} \nl
    &\subseteq \bb{\x[k]} \br{-\x[k]h + \bigo{\pol{\x}h^2}} \nl
    &\subseteq -\x[k]\bb{\x}h +  \bb{\x[k]} \cdot \bigo{\pol{\x}h^2}
  \end{align}
  For any $\seq{\epsilon_n} \in \bigo{\pol{\x[k]}h^2}$ there is a polynomial $p$ and a $A_1 > 0$ with
  \begin{align}
    \abs{\epsilon_n} \le \abs{p(\x[k])} \cdot A_1 h^2
  \end{align}
  The polynomial $p$ is bounded on $I$ so that there is a $A_2>0$ with $\abs{p(x)} \le A_2$ for all $x\in I$. We set $A=A_1\cdot A_2$ and get $\abs{\epsilon_n} \le A h^2$ or $\seq{\epsilon_n}\in \bigabs{A h^2}$. So we can conclude

  \begin{align}
    \bb{\x[k+1]} &\in \bb{\x[k]} -\x[k]\bb{\x[k]} h + \bb{\x} \cdot \bigabs{A h^2}\bnl
    &
    \begin{comment}
      \text{Set } f(x,y) = -xy
    \end{comment} \bnl
    &= \bb{\x[k]} + \br{\f{\x[k], \bb{\x[k]}} + \bb{\x} \cdot \bigabs{A h}}h \bnl
    &
    \begin{comment}
      \text{Set } M(x,y,h,f) = f(x,y) + y \cdot \bigabs{A h}
    \end{comment} \bnl
    &= \bb{\x[k]} + M(\x[k], \bb{\x[k]}, f, h) \cdot h
  \end{align}

  Therefore we can think about $\bb{\x}$ as the results of the one-step method $M(x,y,f,h)=f(x,y)+\bb{\x} \cdot \bigabs{A h}$ whereby we take $\bnp$ for $\xnp$ as the initial value. This one-step method corresponds to the ODE with initial condition

  \begin{align}
    \dphin{x}=f(x,\phin{x})=-x\phin{x},\ \phin{\xnp} = \bb{\xnp}
  \end{align}
  Its solution is (see theorem \ref{chr_ode})
  \begin{align}
    \phin{x} = \bb{\xnp}\fexp{-\frac{x^2}2+\frac{\xnp^2}2} = \frac{\bb{\xnp}}{\fphi{\xnp}} \fphi{x}
  \end{align}
  We find by using the Taylor series
  \begin{align}
    & \phin{x+h} = \phin{x} -x\phin{x}h+\br{\tilde x^2-1}\phin{\tilde x} h^2
  \end{align}
  whereas $\tilde x$ is a number bounded by $x$ and $x+h$. Because the term $\phin{\tilde x} = \frac{\bnp}{\fphi{\xnp}} \fphi{\tilde x}$ and the term $\frac{1}{\fphi{\xnp}} \br{\tilde x^2-1} \fphi{\tilde x}$ is bounded on $I$ there is a $B > 0$ with $\abs{\tilde x^2-1} \phin{\tilde x} h^2 \le \bnp \cdot B h^2$ independent of $\x[k]$. Thus

  \begin{align}
    \frac{\phin{x+h}-\phin{x}}h &= -x\phin{x}+\br{\tilde x^2-1}\phin{x} h \in -x\phin{x} + \bigabs{\bnp \cdot Bh}
  \end{align}
  Now we can estimate the local discretization error
  \begin{align}
    \abs{\tau(\x)} &= \abs{\frac{\phin{\x[k]+h}-\phin{\x}}h -  M(\x,\phin{\x},f,h)} \nl
    &\in \abs{-\x\phin{\x} + \bigabs{\bnp \cdot Bh} +\x\phin{\x} - \bb{\x} \bigabs{Ah}} \nl
    &= \abs{\bigabs{\bnp \cdot Bh}-\bb{\x} \bigabs{Ah}} \nl
    &\subseteq \bigabs{\bnp \cdot Bh}+ \bb{\x} \bigabs{Ah} \nl
    &
    \begin{comment}
      \bb{\x} \le \bnp
    \end{comment} \nl
    &\subseteq \bigabs{\bnp \cdot (B+A)h} \nl
    &
    \begin{comment}
      \text{Set } C=A+B
    \end{comment} \nl
    &= \bigabs{\bnp \cdot Ch}
  \end{align}
  from which follows $\abs{\tau(\x)} \le \bnp \cdot Ch$. Besides approximating $\abs{\tau(x)}$ we also need to prove, that $M(x,y,h,f)$ is Lipschitz continuous regarding $y$. We have
  \begin{align}
    \abs{f(x,y_1)-f(x,y_2)} &= \abs{-xy_1+xy_2} = \abs{x} \cdot \abs{y_1-y_2} \le \Lt \abs{y_1-y_2}
  \end{align}
  and therefore
  \begin{align}
    & M(x,y_1,h,f)-M(x,y_2,h,f) \nl
    \in\ &  f(x,y_1)+y_1 \cdot \bigabs{Ah} - f(x,y_2)-y_2 \cdot \bigabs{A h} \nl
    =\ & f(x,y_1)-f(x,y_2) + y_1 \cdot \bigabs{Ah}-y_2 \cdot \bigabs{Ah} \nl
    \subseteq\ & \br{y_1-y_2} \cdot \bigabs{\Lt} + \br{y_1-y_2} \cdot \bigabs{Ah}  \nl
    \subseteq\ & \br{y_1-y_2} \cdot \bigabs{\Lt+Ah}
  \end{align}
  So we can conclude
  \begin{align}
    \abs{M(x,y_1,h,f)-M(x,y_2,h,f)} \le (\Lt + Ah) \cdot \abs{y_1-y_2}
  \end{align}

  \noindent From the theorem \ref{thm:ode_global_error} about the global error of one-step methods we can conclude for $L=\Lt+Ah$
  \begin{align}
    & \abs{\bb{\x}-\phin{\x}} \le \frac{\fe{L\abs{\x-\xnp}}-1}L \cdot \bnp \cdot Ch \bnl
    \implies & \bb{\x} \in \phin{\x} + \bigabs{\bnp \cdot \frac{\fe{L\abs{\x-\xnp}}-1}L Ch}
  \end{align}
\end{proof}

\section{Global version of the de Moivre-Laplace theorem}

In this section we will prove the statement:
\begin{align}
  \lim_{n\to\infty} \abs{\sum_{a \le \x \le b} \bb{\x[k]}- \int_a^b \fphi{t} \d{t}} = 0
\end{align}
and
\begin{align}
  \lim_{n\to\infty} \abs{\sum_{\x \le z} \bb{\x[k]} - \int_{-\infty}^z \fphi{t} \d{t}} = 0
\end{align}
for any $a,b,z\in\R$. Thereby, we will split the proof in several steps. We start with the following lemma:

\begin{lemma} \label{alt:lemma1}
  For any $a,b \in \R$ with $a \le b$ we have

  \begin{align}
    \lim_{n\to\infty} \abs{\frac{\fphi{\xnp}h}{\bb{\xnp}} \sum_{a\le\x\le b} \bb{\x} - \int_a^b \fphi{t}\d{t}} = 0
  \end{align}
\end{lemma}

\begin{proof}
  From the local version of de Moivre-Laplace with $I=[a,b]$ we can deduce that there is a $C > 0$ independent of $\x$ and $n$ with (see remark \ref{remark:alt_local})
  \begin{align}
    \bb{\x} \in \frac{\bnp}{\fphi{\xnp}}\fphi{\x} + \bigabs{{\bnp}\cdot Ch} \subseteq \frac{\bnp}{\fphi{\xnp}}\fphi{\x} + \bnp \cdot \bigo{h}
  \end{align}
  With this equation follows
  \begin{align}
    \frac{\fphi{\xnp}h}{\bnp} \sum_{a \le \x \le b} \bb{\x} &\in \frac{\fphi{\xnp}h}{\bnp} \sum_{a \le \x \le b}\br{\frac{\bnp}{\fphi{\xnp}}\fphi{\x} + \bnp \cdot \bigo{h}} \bnl
    &
    \begin{comment}
      \lim_{n\to\infty} \fphi{\xnp} = \frac{1}{\sqrt{2\pi}} \implies \seq{\fphi{\xnp}} \in \bigo{1}
    \end{comment} \bnl
    &\subseteq \sum_{a \le \x \le b}\br{\fphi{\x}h + \bigo{h^2}}\bnl
    &
    \begin{comment}
      \begin{aligned}
        \sum_{a \le \x \le b} \fphi{\x}h &\in \int_a^b \fphi{t} \d{t} + \bigo{h} \bnl
        \sum_{a \le \x \le b} \bigo{h^2} &\subseteq \bigo{h}
      \end{aligned}
    \end{comment} \bnl
    &\subseteq \int_a^b \fphi{t} \d{t} + \bigo{h}
  \end{align}
\end{proof}

\begin{lemma} \label{lemma_ineq}
  All $x\in\R$ with $0 < x \le \frac 12$ fulfill the inequality $\frac{1+x}{1-x} \le 1+4x$.
\end{lemma}

\begin{proof}
  \begin{align}
    \begin{array}{rrl}
      & x & \le \frac 12 \nl
      \implies & 4x^2 & \le 2x \nl
      \implies & 0 & \le  2x - 4x^2 \nl
      \implies & 1+x & \le  1+3x-4x^2 = (1+4x)(1-x) \nl
      &&
      \begin{comment}
        1-x > 0
      \end{comment} \nl
      \implies & \frac{1+x}{1-x} & \le 1+4x
    \end{array}
  \end{align}
\end{proof}

\begin{lemma}
  We have $\lim_{n\to\infty} \frac{\fphi{\xnp}h}{\bnp} = 1$.
\end{lemma}

\begin{remark}
  This lemma implies $\bnp \sim \frac{1}{\sqrt{2\pi npq}}$ because $\lim_{n\to\infty} \fphi{\xnp} = \frac1{\sqrt{2\pi}}$ as we already know from the standard proof of the de Moivre-Laplace theorem based on Stirling's formula. Therefore this lemma can be shown with Stirling's formula. Yet it is only applicable for the binomial distribution and a use of Stirling's formula at this point contradicts the attempt to find a new proof of de Moivre-Laplace (which does not depend on methods of the standard proof). Thus I want to omit this approximation.
\end{remark}

\begin{proof}
  Let $c_n = \frac{\fphi{\xnp}h}{\bnp}$. Let $\epsilon > 0$ be arbitrary with $\epsilon \le 2$. From the Chebyshev's inequality we know $\P{\abs{X}> R} \le \frac{1}{R^2}$ for any $R > 0$ and any standardized random variable $X$ \cite{wiki:chebyshev}. Thus there is an $R >0$ such that
  \begin{align}\label{eq1}
    1 -\frac{\epsilon}4 \le \sum_{-R \le \x \le R} \bb{\x[k]} \le 1
  \end{align}
  and
  \begin{align}\label{eq2}
    1-\frac{\epsilon}4 \le \int_{-R}^R \fphi{t}\d{t} \le 1
  \end{align}

  From the above lemma \ref{alt:lemma1} we know that there is a $N\in\N$ so that for each $n\in\N$ with $n\ge N$ we find
  \begin{align} \label{eq3}
    \abs{c_n \sum_{-R\le\x\le R}\bb{\x} - \int_{-R}^R \fphi{t}\d{t}} \le \frac{\epsilon}4
  \end{align}
  The inequalities \eqref{eq2} and \eqref{eq3} imply
  \begin{align}
    1-\frac{\epsilon}2 \le c_n \sum_{-R \le \x \le R} \bb{\x} \le 1+\frac{\epsilon}4
  \end{align}
  Together with the first inequality \eqref{eq1} we can conclude
  \begin{align} \label{shown1}
    1-\frac \epsilon2 \le c_n \sum_{-R \le \x \le R} \bb{\x} \le c_n
  \end{align}
  and
  \begin{align}
    & 1+\frac \epsilon4 \ge c_n \sum_{-R \le \x \le R} \bb{\x} \ge c_n \br{1-\frac \epsilon4} \bnl
    \implies & c_n \le \frac{1+\frac \epsilon4}{1-\frac\epsilon4} \bnl
    &
    \begin{comment}
      \text{Lemma \ref{lemma_ineq}}
    \end{comment} \bnl
    \implies & c_n \le 1+4\cdot \frac \epsilon 4 = 1+\epsilon \label{shown2}
  \end{align}
  So we get from \eqref{shown1} and \eqref{shown2}
  \begin{align}
    1-\frac \epsilon 2 \le c_n \le 1+\epsilon
  \end{align}
  which proves $\abs{c_n-1} \le \epsilon$ and thus $\lim_{n\to\infty} c_n = 1$.
\end{proof}

\begin{theorem}[First global version of the de Moivre-Laplace theorem]
  For any $a,b\in\R$ we find

  \begin{align}
    \lim_{n\to\infty} \abs{\sum_{a \le \x \le b} \bb{\x} - \int_{a \le t \le b} \fphi{t} \d{t} } = 0
  \end{align}
\end{theorem}

\begin{proof}
  For $a > b$ the theorem is trivial as
  \begin{align}
    \sum_{a \le \x \le b} \bb{\x} = \int_{a \le t \le b} \fphi{t} \d{t} = 0
  \end{align}
  So we consider only the case $a\le b$. In lemma \ref{alt:lemma1} we have already shown that
  \begin{align}
    \frac{\fphi{\xnp}h}{\bnp} \sum_{a\le \x \le b} \bb{\x} \in \int_a^b \fphi{t} \d{t} + \bigo{h}
  \end{align}
  Therefore
  \begin{align}
    \sum_{a\le \x \le b} \bb{\x} & \in \frac{\bnp}{\fphi{\xnp}h} \br{\int_a^b \fphi{t} \d{t} + \bigo{h}} \bnl
    &
    \begin{comment}
      \lim_{n\to\infty} \frac{\bnp}{\fphi{\xnp}h} = 1 \implies \frac{\bnp}{\fphi{\xnp}h} = 1+\littleo{1}
    \end{comment} \bnl
    & \subseteq \br{1+\littleo{1}}\cdot \br{\int_a^b \fphi{t} \d{t} + \bigo{h}} \bnl
    & \subseteq \int_a^b \fphi{t}\d{t} + \littleo{1} \cdot \int_a^b \fphi{t} \d{t} + \bigo{h} + \littleo{1} \cdot \bigo{h} \bnl
    &
    \begin{comment}
      \begin{aligned}
        \int_a^b \fphi{t} \d{t} & \le 1 \nl
        \bigo{h} & \subseteq \littleo{1}
      \end{aligned}
    \end{comment} \bnl
    & \subseteq \int_a^b \fphi{t} \d{t} + \littleo{1}
  \end{align}
\end{proof}

\begin{theorem}[Second global version of the de Moivre-Laplace theorem]
  For any $z\in\R$ we have

  \begin{align}
    \lim_{n\to\infty} \abs{\sum_{\x \le z} \bb{\x} - \int_{-\infty}^z \fphi{t}\d{t} } = 0
  \end{align}
\end{theorem}

\begin{proof}
  For any $c \le z$ we have the estimation

  \begin{align}
    \abs{\sum_{\x \le z} \bb{\x} - \int_{-\infty}^z \fphi{t}\d{t} } \le
    \underbrace{\abs{\sum_{\x \le c} \bb{\x} - \int_{-\infty}^c \fphi{t}\d{t} }}_\text{term I} + \underbrace{\abs{\sum_{c\le \x \le z} \bb{\x} - \int_c^z \fphi{t}\d{t} }}_\text{term II}
  \end{align}

  First we take $c$ small enough such that term I is smaller than $\frac \epsilon 2$. Due to the Chebyshev's inequality \cite{wiki:chebyshev} we know that such a $c$ must exist. From the first global version of de Moivre-Laplace we know that there is a $N \in\N$ such that the term II is smaller then $\frac \epsilon 2$ for all $n\ge N$. Thus we can make the above sum of the two terms arbitrary small for large enough $n\in\N$ which proves the convergence.
\end{proof}

  \chapter{Conclusion}

In this thesis I proposed a new proof for the de Moivre-Laplace theorem, which could represent a valid alternative to the standard proof concerning error estimation and concerning comprehensibility in class or in textbooks. 

However, compared with the classic method, using Stirling’s formula, my alternative proof could not yield better approximations or smaller error bounds. This was already prevalent when testing it against the local version of the de Moivre-Laplace theorem.

However the new proof has some clear advantages over the standard one:

\begin{enumerate}
  \item My new proof only uses the following property of the binomial distribution
    \begin{align}
      \bb{\x[k+1]} \in \bb{\x[k]} - \x[k]\bb{\x[k]} h + \bigo{\bb{\x[k]} \pol{\x[k]} h^2}
    \end{align}
    Therefore, the same proof applies to each discrete probability distribution fulfilling this property.

  \item My alternative proof connects the theory of ordinary differential equations in an interesting way with probability theory: The one-step method corresponds to the given function and the solution of the ODE is the approximation. Thus, the role of approximated function and approximation is exchanged.

  \item Eventhough, it is a very subjective decision to rank one proof more understandable than another \cite{tampis:understandability}, the new proof arguably is easy to understand and intuitive. Therefore, it could be used in courses or textbooks for undergraduates. An investigation on this issue is out of scope for this thesis.
\end{enumerate}

Eventhough my alternative proof does not result in better approximations or smaller error bounds, the presented proof is a valuable alternative for teaching and understanding the de Moivre-Laplace theorem and an addition to the classic method. 


  \begin{appendices}
    \chapter{The Berry-Esseen theorem}

In this chapter I present you the Berry-Essen theorem which not only proves the central limit but also gives an upper bound for the maximal error of the normal approximation. In the proof of this theorem I will mainly follow Nourdin and Peccati\cite[p. 71 ff.]{nourdin}\todo{Bessere Referenz - Originalquelle?}, who use the Method of Stein.

\section{CLT and Berry-Esseen bounds}

From the central limit theorem we know, that each standardized sum $S_n = \tfrac 1{\sqrt n\sigma} \left(\sum_{k=1}^n x_k - n\mu\right)$ of i.i.d real valued random variables $\seq{X_n}$ with finite expected value $\mu=\E{X_1}$ and finite variance $\sigma^2=\E{X_1^2}-\E{X_1}^2$ converges in distribution to the standard normal distribution \refneeded. So we have $\lim_{n\to\infty} \P{S_n\le x} = \Phi(x)$ for all $x\in\R$. Thus we can use the normal distribution to approximate the cumulative distribution function of $S_n$.

Unfortunately not each proof of the central limit theorem provides an estimate of the error $\abs{\P{S_n\le x}-\Phi(x)}$ \refneeded. Here comes the Berry-Esseen theorem into play. It provides an upper bound for the maximal error between $\P{S_n \le x}$ and $\Phi(x)$ \todo{$\Phi$ muss eingeführt werden}. Therefore the Berry-Esseen theorem needs the finiteness of the third moment for each $X_n$ as an additional premise compared to the central limit theorem. The Berry-Esseen theorem states\cite[p. 71]{nourdin}:

\begin{theorem}[CLT and Berry-Esseen bounds]
  Let $\seq{X_n}$ be a sequence of standardized i.i.d real valued random variables with finite third moment. Thus we have $\E{X_1}=0$, $\E{X_1^2}=1$ and $\E{|X_1|^3} < \infty$. Let $S_n$ be the associated sequence of normalized partial sums, i.e

  \begin{align}
    S_n = \frac{1}{\sqrt n} \sum_{k=1}^n X_k
  \end{align}

  The partial sums $S_n$ converge in law to the standard normal distribution. Moreover there is a $C > 0$ not depending on $n$ nor $\seq{X_n}$ such that

  \begin{align}
    \sup_{x\in\R} \abs{\P{S_n\le x}-\Phi(x)} \le \frac{C \E{|X_1|^3}}{\sqrt n}
  \end{align}
\end{theorem}

\begin{remark}
  The currently best known estimate for the best choice of $C$ is $C \ge \Cmin$\cite{esseen1956} and $C \le \Cmax$\cite{shevtsova2011}.
\end{remark}

\begin{remark}[Berry-Esseen theorem for the binomial distribution]
  If for example $\seq{X_n}$ are binomial trials with the probability of success $p \in (0,1)$ the third moment of all $X_n$ are finite. With $q=1-p$ we get \refneeded

\begin{align}
  \E{X_1}=\frac{p^2+q^2}{\sqrt{pq}}
\end{align}

Thus the Berry-Esseen theorem yields the following estimate for the normalized binomial distribution

\begin{align}
  \sup_{x\in\R} \abs{\P{S_n\le x}-\Phi(x)} \le \frac{C(p^2+q^2)}{\sqrt{npq}}
\end{align}

\noindent Thereby the best choice of $C$ fulfills $C \le \Cmax$.
\end{remark}

\section{The method of Stein}

\subsection{The basic idea of the method of Stein}

We need to estimate $\sup_{x\in\R} \abs{\P{S_n\le x}-\Phi(x)}$ to prove the Berry-Esseen theorem. To do so we first rewrite this supremum norm:

\begin{align}
   & \sup_{x\in\R} \abs{\P{S_n\le x}-\Phi(x)} \nl
   &
   \begin{comment}
     \text{Let } N \text{ have a standard normal distribution}
   \end{comment} \nl
  =& \sup_{x\in\R} \abs{\int_{-\infty}^x \d{S_n} - \int_{-\infty}^x \d{N}} \nl
  =& \sup_{x\in\R} \abs{\int_{\R} \cfs{x} \d{S_n} - \int_{\R} \cfs{x} \d{N}} \nl
  =& \sup_{x\in\R} \abs{\E{\cf{x}{S_n}} - \E{\cf{x}{N}}} \nl
   &
   \begin{comment}
     \text{Let }\mathcal H=\left\{\cfs{x} : x\in\R\right\}
   \end{comment} \nl
  =& \sup_{h\in\mathcal H} \abs{\E{h(S_n)}-\E{h(N)}}
\end{align}

Now we can use the method of Stein to transform the difference $\E{h(S_n)}-\E{h(N)}$ in a way that $N$ does not occur any more. Here we take the solution $f$ of the following so called \emph{Stein's equation}:

\begin{align}
  f'(x)=xf(x) + h(x) - \E{h(N)}
\end{align}

Let's assume that this ordinary differential equation has a solution $f$. We get for any random variable $Y$:

\begin{align}
  &            & f'(x) - xf(x) & = h(x) - \E{h(N)} \nl
  &  \implies\ & \int_{\R} (f'(t)-t f(t)) \d{Y}(t) &= \int_{\R} (h(t)-\E{h(N)}) \d{Y}(t) \nl 
  &  \implies\ & \E{f'(Y)-Y f(Y)} &= \E{h(Y)}-\E{h(N)}
\end{align}

So instead of $\abs{\E{h(S_n)}-\E{h(N)}}$ we can estimate $\abs{\E{f'(S_n)-S_n f(S_n)}}$ which does not contain $N$ anymore. Thereby we will use some properties of $f$ which follow from the Stein's equation, namely $\snorm{f} \le \sqrt{\tfrac \pi2}$ and $\snorm{f'} \le 2$.

Before I continue with the proof of the Berry-Esseen theorem I want to concrete what I have mentioned in this section.

\subsection{Stein's equation and its solutions}

For this and the following section we assume that $N$ is a random variable with a standard normal distribution. We have

\begin{definition}[Stein's equation]
  Let $h:\R \to \R$ be a Borel function such that $\E{h(N)}$ exists and $\abs{\E{h(N)}} < \infty$. The \emph{Stein's equation} associated to $h$ is the following ordinary differential equation

  \begin{align}
    f'(x) = x f(x) + h(x) - \E{h(N)}
  \end{align}
\end{definition}

The following proposition shows that the Stein's equation always has a solution and provides an explicit form for it:

\begin{proposition}[Solutions for the Stein's equation]
  Let $f_h$ be defined by

  \begin{align}
    f_h(x) = \e{\frac{x^2}{2}} \int_{-\infty}^x \Big[h(y)-\E{h(N)}\Big] \e{-\frac{y^2}{2}}\d{y}
  \end{align}

  \noindent Every solution to the Stein's equation has the form

  \begin{align}
    f(x) = c \e{\frac{x^2}{2}} + f_h(x)
  \end{align}

  \noindent with $c \in \R$.
\end{proposition}

\begin{proof}
   Because $\E{h(N)}$ exists \todo{bessere Formulierung}, the function $y\mapsto h(y)\e{-\frac{y^2}{2}}$ is integrable so that also for each $x\in\R$ the integral ${\int_{-\infty}^x \Big[h(y)-\E{h(N)}\Big]\e{-\frac{y^2}{2}}\d{y}}$ exists. Thus $f_h$ is well defined. From the Stein's equation follows

  \begin{align}
    &      & f'(x)-xf(x) & = h(x) - \E{h(N)} \nl
    & \iff & \left[ f'(x) - xf(x) \right] \e{-\frac{x^2}{2}} &= \Big[h(x) - \E{h(N)}\Big] \e{-\frac{x^2}{2}} \nl
    & \iff & \partial_x \left( \e{-\frac{x^2}2} f(x) \right) &= \Big[h(x) - \E{h(N)}\Big] \e{-\frac{x^2}{2}} \nl
    & \iff & \e{-\frac{x^2}{2}} f(x) & = \int_{-\infty}^x \Big[ h(y) - \E{h(N)} \Big] \e{-\frac{y^2}{2}} \d{y} + c \nl
    & \iff & f(x) & = f_h(x) + c \e{\frac{x^2}{2}}
  \end{align}
\end{proof}

In the proof of the Berry-Esseen theorem we will need a estimate of $\snorm{f_h}$ and $\snorm{f_h'}$. Therefore we need to restrict the values of $h$ in the interval $h$. After this restriction we get

\begin{proposition}[Stein's bounds]
  Let $h:\R \to [0,1]$ be a Borel-function such that $\E{h(N)}$ exists with $\E{h(N)}<\infty$. The solution $f_h$ of Stein's equation given in the above proposition fulfills

  \begin{align}
    \snorm{f_h} \le \sqrt{\frac \pi2} \text{ and } \snorm{f_h'} \le 2
  \end{align}
\end{proposition}

\begin{proof}
  We start with proving $\snorm{f_h}\le \sqrt{\frac{\pi}{2}}$. First get the estimate

  \begin{align} \label{eq:estimate1}
    \abs{f_h(x)} & \le \abs{\e{-\frac{x^2}{2}} \int_{-\infty}^x \Big[ h(y) - \E{h(N)} \Big] \e{-\frac{y^2}{2}} \d{y} } \nl 
    &\le \e{\frac{x^2}{2}} \int_{-\infty}^x \abs{h(y)-\E{h(N)}} \e{-\frac{y^2}{2}}\d{y} \nl
    &
    \begin{comment}
      h(y), \E{h(N)} \in [0,1] \implies \abs{h(y)-\E{h(N)}} \le 1
    \end{comment} \nl
    & \le \e{\frac{x^2}{2}} \int_{-\infty}^x \e{-\frac{y^2}{2}} \d{y}
  \end{align}
  
  \noindent We also have

  \begin{align}
    \int_{-\infty}^x \Big[h(y)-\E{h(N)}\Big] \e{-\frac{y^2}{2}} \d{y} = -\int_{x}^\infty \Big[h(y)-\E{h(N)}\Big] \e{-\frac{y^2}{2}} \d{y}
  \end{align}

  \noindent because $\int_{\R} \Big[h(y)-\E{h(N)}\Big] \e{-\frac{y^2}{2}} \d{y}=0$ \todo{$\d{N(y)}$ richtig?}:

  \begin{align} \label{eq:estimate2}
    \int_{\R} \Big[ h(y)-\E{h(N)} \Big] \e{-\frac{y^2}{2}} \d{y} & = \int_{\R} h(y) \e{-\frac{y^2}{2}} \d{y} - \int_{\R} \E{h(N)} \e{-\frac{y^2}2} \d{y} \nl
    &= \sqrt{2\pi} \int_{\R} h(y) \frac{1}{2\pi} \e{-\frac{y^2}{2}} \d{y} - \sqrt{2\pi} \E{h(N)} \int_{\R} \frac{1}{\sqrt{2\pi}} \e{-\frac{y^2}{2}} \d{y} \nl
    &= \sqrt{2\pi} \int_{\R} h(y) \d{N(y)} - \sqrt{2\pi} \E{h(N)} \int_{\R} \d{N(y)} \nl
    &= \sqrt{2\pi}\E{h(N)} - \sqrt{2\pi}\E{h(N)} \nl
    &= 0
  \end{align}

  \noindent Thus we also could estimate $\abs{f_h(x)}$ via

  \begin{align}
    \abs{f_h(x)} & \le \abs{\e{-\frac{x^2}{2}} \int_{-\infty}^x \Big[ h(y) - \E{h(N)} \Big] \e{-\frac{y^2}{2}} \d{y} } \nl 
    & \le \abs{-\e{-\frac{x^2}{2}} \int_{x}^\infty \Big[ h(y) - \E{h(N)} \Big] \e{-\frac{y^2}{2}} \d{y} } \nl 
    &\le \e{\frac{x^2}{2}} \int_{x}^\infty \abs{h(y)-\E{h(N)}} \e{-\frac{y^2}{2}}\d{y} \nl
    &\le \e{\frac{x^2}{2}} \int_{x}^\infty \e{-\frac{y^2}{2}} \d{y}
  \end{align}
 
  \noindent From \eqref{eq:estimate1} and \eqref{eq:estimate2} follows

  \begin{align}
    \abs{f_h(x)} & \le \min\left\{ \e{\frac{x^2}{2}} \int_{-\infty}^x \e{-\frac{y^2}{2}} \d{y}, \e{\frac{x^2}{2}} \int_x^\infty \e{-\frac{y^2}{2}} \d{y} \right\} \nl
    & \le \e{\frac{x^2}{2}} \min\left\{ \int_{-\infty}^x \e{-\frac{y^2}{2}} \d{y}, \int_x^\infty \e{-\frac{y^2}{2}} \d{y} \right\} \nl
    &
    \begin{comment}
      y \mapsto \e{-\frac{y^2}{2}} \text{ is even}
    \end{comment} \nl
    & = \underbrace{\e{\frac{x^2}{2}} \int_{\abs{x}}^\infty \e{-\frac{y^2}{2}} \d{y}}_{=: s(x)} \nl
  \end{align}

  \noindent To find the maximum of $s(x)$ we first compute $s'(x)$ for $x > 0$:

  \begin{align}
    s'(x) &= \partial_x \left( \e{\frac{x^2}{2}} \int_x^\infty \e{-\frac{y^2}{2}} \d{y} \right) \nl
    &= x\e{\frac{x^2}{2}} \int_x^\infty \e{-\frac{y^2}{2}} \d{y} - 1 \nl
    &= \e{\frac{x^2}{2}} \int_x^\infty x\e{-\frac{y^2}{2}} \d{y} - 1 \nl
    &
    \begin{comment}
      x \le y \text{ on } [x,\infty)
    \end{comment} \nl
    &\le \e{\frac{x^2}{2}} \int_x^\infty y\e{-\frac{y^2}{2}} \d{y} - 1 \nl
    &= \e{\frac{x^2}{2}} \left[-\e{-\frac{y^2}{2}}\right]_x^\infty \d{y} - 1 \nl
    &= \e{\frac{x^2}{2}} \e{-\frac{x^2}{2}} - 1 \nl
    &= 0
  \end{align}

  Thus $s'(x) \le 0$ for $x \ge 0$ so that $s(x)$ is monotone decreasing on $[x,\infty)$. Because $s(x)$ is even (i.e. $s(-x)=s(x)$) we have $s'(x) \ge 0$ for $x < 0$ which implies that $s(x)$ is monotone increasing on $\R^{-}_0$. Altogether $s(x)$ has its global maximum at $x=0$ which is

  \begin{align}
    s(0) = \int_0^\infty \e{-\frac{y^2}{2}} \d{y} = \frac{1}{2} \int_{-\infty}^\infty \e{-\frac{y^2}{2}} \d{y} = \frac{1}{2} \sqrt{2\pi} = \sqrt{\frac{\pi}{2}}
  \end{align}

  \noindent Finally we have

  \begin{align}
    \abs{f_h(x)} \le s(x) \le \sqrt{\frac \pi 2}
  \end{align}

  \noindent To prove the second estimate $\snorm{f_h'} \le 2$ we can start with Stein's equation:

  \begin{align}
    \abs{f_h'(x)} & = \abs{x f_h(x) + h(x) - \E{h(N)}} \nl
    & = \abs{x \e{\frac{x^2}{2}} \int_{-\infty}^x \Big[ h(y) - \E{h(N)} \Big] \e{-\frac{y^2}{2}} \d{y}} \nl
    & \le \abs{x} \e{\frac{x^2}{2}} \int_{-\infty}^x \abs{h(y)-\E{h(N)}} \e{-\frac{y^2}{2}} \d{y}  + \abs{h(y)-\E{h(y)}} \nl
    &
    \begin{comment}
      \abs{h(y)-\E{h(N)}} \le 1
    \end{comment} \nl
    & \le \abs{x} \e{\frac{x^2}{2}} \int_{-\infty}^x \e{-\frac{y^2}{2}} \d{y}  + 1 \nl
  \end{align}

  \noindent Again we can use the following equation in the estimation of $\abs{f_h'(x)}$

  \begin{align}
    \int_{-\infty}^x \Big[h(y)-\E{h(N)}\Big] \e{-\frac{y^2}{2}} \d{y} = -\int_{x}^\infty \Big[h(y)-\E{h(N)}\Big] \e{-\frac{y^2}{2}} \d{y}
  \end{align}

  \noindent No we get

  \begin{align}
    \abs{f_h'(x)} \le \abs{x} \e{\frac{x^2}{2}} \int_x^\infty \e{-\frac{y^2}{2}} \d{y}  + 1 \nl
  \end{align}

  \noindent Altogether we have

  \begin{align}
    \abs{f_h'(x)} & \le \abs{x} \e{\frac{x^2}{2}} \min\left\{ \int_{-\infty}^x \e{-\frac{y^2}{2}} \d{y},  \int_x^\infty \e{-\frac{y^2}{2}} \d{y} \right\} + 1 \nl
    & = \abs{x} \e{\frac{x^2}{2}} \int_{\abs{x}}^\infty \e{-\frac{y^2}{2}} \d{y} + 1 \nl
    & = \e{\frac{x^2}{2}} \int_{\abs{x}}^\infty \abs{x} \e{-\frac{y^2}{2}} \d{y} + 1 \nl
    &
    \begin{comment}
      \abs{x} \le y \text{ on interval } \left[\abs{x},\infty\right)
    \end{comment} \nl
    & = \e{\frac{x^2}{2}} \int_{\abs{x}}^\infty y \e{-\frac{y^2}{2}} \d{y} + 1 \nl
    & = \e{\frac{x^2}{2}} \e{-\frac{x^2}{2}} + 1 \nl
    & = 2 \nl
  \end{align}

  \noindent This proves $\snorm{f_h'} \le 2$.
\end{proof}

\section{Proof of the Berry-Esseen theorem}

\subsection{First steps of the proof\todo{better title}}

Let's recall the Berry-Esseen theorem: For the standardized partial sum $S_n$ of i.i.d random variables $\seq{X_n}$ with finite third moment $\E{\abs{X_n}^3}<\infty$ there is a $C > 0$ such that

\begin{align}
  \sup_{x\in\R} \abs{\P{S_n\le x}-\Phi(x)} \le \frac{C \E{|X_1|^3}}{\sqrt n}
\end{align}

\noindent We first define $C_n$ via

\begin{align}
  C_n = \frac{\sqrt{n}}{\E{\abs{X_1}^3}} \sup_{x\in\R} \abs{\P{S_n \le x}-\Phi(x)}
\end{align}

\noindent Thus $C_n$ is the smallest real number with

\begin{align}
  \sup_{x\in\R}\abs{\P{S_n\le x} - \Phi(x)} \le \frac{C_n\E{\abs{X_1}^3}}{\sqrt{n}}
\end{align}

To prove the Berry-Essen theorem we have to show, that $\sup_{n\in\N} C_n < \infty$. Let $\mathcal H$ be the set of all $\cfs{x}$ for $x\in\R$. We already have shown in the section about Stein's method that

\begin{align}
  \sup_{x\in\R}\abs{\P{S_n\le x} - \Phi(x)} & = \sup_{h\in\mathcal H} \abs{\E{h(S_n)}-\E{h(N)}} \nl
  &= \sup_{h\in\R} \abs{\E{f_h'(S_n)-S_n f_h(S_n)}}
\end{align}

Unfortunately the functions of $\mathcal H$ are not continuous and thus $f_h$ is no $C^1$-function for any $h\in \mathcal H$. That's why we introduce the functions $h_{\epsilon,x}$ which are defined by

\includefig{continuous_step_function}{The function $h_{\epsilon,x}$}{The function $h_{epsilon,x}$}

\noindent For any random variable $Y$ it is

\begin{align} \label{berryproof:inequality1}
  \E{h_{x-\epsilon,\epsilon}(Y)} \le \P{Y \le x} \le \E{h_{x+\epsilon,\epsilon}(Y)}
\end{align}

Because the density function $\phi(t)$ of $N$ with the standard normal distribution is bounded by $\tfrac{1}{\sqrt{2\pi}}$ we have:

\begin{align}
  \E{h_{x+\epsilon,\epsilon}(N)}-\frac{4\epsilon}{\sqrt{2\pi}} \le \E{h_{x-\epsilon,\epsilon}(N)} \text{ and } \E{h_{x+\epsilon,\epsilon}(N)} \le \E{h_{x-\epsilon,\epsilon}(N)} + \frac{4\epsilon}{\sqrt{2\pi}}
\end{align}

\noindent Together with \eqref{berryproof:inequality1} we get

\begin{align} \label{berryproof:inequality2}
  \E{h_{x+\epsilon,\epsilon}(N)}-\frac{4\epsilon}{\sqrt{2\pi}} \le \P{N\le x} \le \E{h_{x-\epsilon,\epsilon}(N)} + \frac{4\epsilon}{\sqrt{2\pi}}
\end{align}

\noindent When we combine \eqref{berryproof:inequality1} and \eqref{berryproof:inequality2} we get the estimate

\begin{align}
  \E{h_{x-\epsilon,\epsilon}(S_n)} - \E{h_{x-\epsilon,\epsilon}(N)} - \frac{4\epsilon}{\sqrt{2\pi}} & \le \P{S_n \le x} - \P{N \le x} \nl
  &\le \E{h_{x+\epsilon,\epsilon}(S_n)} - \E{h_{x+\epsilon,\epsilon}(N)} + \frac{4\epsilon}{\sqrt{2\pi}}
\end{align}

\noindent Thus

\begin{align}
  \abs{\P{S_n \le x}-\P{N \le x}} \le \max\left\{
    \begin{array}{l}
      \E{h_{x-\epsilon,\epsilon}(S_n)} - \E{h_{x-\epsilon,\epsilon}(N)}, \nl
      \E{h_{x+\epsilon,\epsilon}(S_n)} - \E{h_{x+\epsilon,\epsilon}(N)}
    \end{array}
  \right\} + \frac{4\epsilon}{\sqrt{2\pi}}
\end{align}

\noindent So we get the estimate

\begin{align} \label{berryproof:before_stein}
  \sup_{x\in\R} \abs{\P{S_n \le x} - \P{N \le x}} \le \sup_{x\in\R} \Big\{ \abs{\E{h_{x,\epsilon}(S_n)} - \E{h_{x,\epsilon}(N)}} \Big\} + \frac{4\epsilon}{\sqrt{2\pi}}
\end{align}

For estimating $\abs{\E{h_{x,\epsilon}(S_n)} - \E{h_{x,\epsilon}(N)}}$ we can now use the method of Stein. Because $h_{x,\epsilon}$ is continuous its solution to the Stein's equation is a $C^1$-function.

\subsection{Application of the method of Stein}

In the following I set $h=h_{x,\epsilon}$. Let $f=f_h$ be the solution to the Stein's equation for $h=h_{x,\epsilon}$. Besides let $n \ge 2$ and let $\tilde f(x)=xf(x)$. I also define $S_n^k$ via

\begin{align}
  S_n^k = \frac 1{\sqrt{n}} \sum_{j\neq k} Y_j = S_n - \frac 1{\sqrt n} Y_k
\end{align}

\noindent In the section about the method of Stein we have seen

\begin{align}
  \E{h(S_n)}-\E{h(N)}=\E{f'(S_n)-S_n f(S_n)}
\end{align}

\noindent The later term can be rewritten as

\begin{align}
  \E{h(S_n)}-\E{h(N)} &= \E{ \fp{S_n} -S_nf(S_h)} \nl
  & \begin{comment}
    S_n = \frac 1{\sqrt n} \sum_{k=1}^n Y_k
  \end{comment} \nl
  &= \E{f'(S_n)-\sum_{k=1}^n \frac{Y_k}{\sqrt n} f(S_n)} \nl
  &= \sum_{k=1}^n \E{f'(S_n)\frac 1n-\frac{Y_k}{\sqrt n} f(S_n)} \nl
  & \begin{comment}
    \E{ \f{S_n^k} Y_k } = \E{ \f{S_n^k} }\E{Y_k} = 0
  \end{comment} \nl
  &= \sum_{k=1}^n \E{ f'(S_n) \frac 1n- \frac{Y_k}{\sqrt n} f(S_n) + \frac{Y_k}{\sqrt n} \f{S_n^k }} \nl
  &= \sum_{k=1}^n \E{ f'(S_n) \frac 1n - \frac{Y_k}{\sqrt n} \left( f(S_n)-\f{S_n^k} \right)} \nl
  &
  \begin{comment}
    \begin{aligned}
  f(y)-f(x) &= \int_x^y f'(s) \d{s} \nl
  &= (y-x) \int_0^1 f'(x+(y-x)s) \d{s} \nl
  &
  \begin{comment}
    \text{\parbox{0.4\textwidth}{
      Choose $\Theta$ uniformly distributed in $[0,1]$ and independent of all $Y_1,\ldots,Y_k$
    }}
  \end{comment} \nl
  &= (y-x)\ \E{ \fp{x+(y-x)} \Theta }
    \end{aligned}
  \end{comment} \nl
  &= \sum_{k=1}^n \E{ f'(S_n) \frac 1n - \frac{Y_k^2}{n} \E{ \fp{ S_n^k + \Theta\frac{Y_k}{\sqrt n} }}} \nl
  &
  \begin{comment}
    \begin{aligned}
      \E{ \frac{Y_k^2}{n} \E{ \fp{ S_n^k + \Theta\frac{Y_k}{\sqrt n} }}} &= \frac{\E{Y_k^2}}{n}  \E{ \ft{S_n^k+\Theta\frac{Y_k}{\sqrt n} }} \nl
&= \frac{1}{n} E\left[f'\left(S_n^k+\Theta\frac{Y_k}{\sqrt n}\right)\right]
    \end{aligned}
  \end{comment} \nl
  &= \frac 1n \sum_{k=1}^n \E{ f'(S_n) - \fp{S_n^k+\Theta\frac{Y_k}{\sqrt n}}} \nl
  &\begin{comment}
    f'(x)=\tilde f(x) + h(x)-E[h(N)]
  \end{comment} \nl
  &= \frac 1n \sum_{k=1}^n \E{
    \begin{multlined}
      \tilde f(S_n) + h(S_n) - E[h(N)] \nl
      -\ft{S_n^k +\Theta \frac{Y_k}{\sqrt n} } - \h{ S_n^k +\Theta \frac{Y_k}{\sqrt n} } + E[h(N)]
    \end{multlined}
  } \nl
  &= \frac 1n \sum_{k=1}^n \E{
  \begin{multlined}
      \tilde f(S_n) + h(S_n) \nl
      -\ft{ S_n^k +\Theta \frac{Y_k}{\sqrt n}} - \h{S_n^k +\Theta \frac{Y_k}{\sqrt n}}
  \end{multlined}
  } \nl
  &= \frac 1n \sum_{k=1}^n \E{
    \begin{multlined}
      \ft{S_n} - \ft{S_n^k} \nl
      - \ft{ S_n^k +\Theta \frac{Y_k}{\sqrt n} } + \ft{ S_n^k } \nl
      + \h{S_n} - \h{ S_n^k } \nl
      - \h{ S_n^k +\Theta \frac{Y_k}{\sqrt n} } + \h{ S_n^k }
    \end{multlined}
  }
\end{align}

\noindent Thus

\begin{align}
  \abs{\E{h(S_n)}-\E{h(N)}}
  &= \abs{\frac 1n \sum_{k=1}^n \E{
    \begin{multlined}
      \ft{S_n} - \ft{S_n^k} \nl
      - \ft{ S_n^k +\Theta \frac{Y_k}{\sqrt n} } + \ft{ S_n^k } \nl
      + \h{S_n} - \h{ S_n^k } \nl
      - \h{ S_n^k +\Theta \frac{Y_k}{\sqrt n} } + \h{ S_n^k }
    \end{multlined}
  }} \nl
  &\le \frac 1n \sum_{k=1}^n \left(
  \begin{multlined}
      \E{\abs{\ft{S_n} - \ft{S_n^k}}} \nl
      +\E{\abs{\ft{ S_n^k +\Theta \frac{Y_k}{\sqrt n} } - \ft{ S_n^k }}} \nl
      +\E{\abs{\h{S_n} - \h{ S_n^k }}} \nl
      +\E{\abs{\h{ S_n^k +\Theta \frac{Y_k}{\sqrt n} } - \h{ S_n^k }}}
    \end{multlined}
  \right)
\end{align}

\noindent So we have four terms which we will estimate respectively

\subsection{Estimation of all four terms}

\subsubsection{Estimation of the first term}

To estimate the first term we will use the fact, that $\snorm{f} \le \sqrt{\tfrac \pi2}$ and $\snorm{f'} \le 2$:

\begin{align}
  &\E{\abs{ \ft{S_n} - \ft{S_n^k} }} \nl
  &
  \begin{comment}
    \begin{aligned}
      |\tilde f(x) - \tilde f(y)| & = |xf(x)-yf(y)| \nl
      & =|xf(x)-yf(x)+yf(x)-yf(y)| \nl
      & =|f(x)(x-y) + y(f(x)-f(y))| \nl
      & \le |f(x)||x-y|+|y||f(x)-f(y)| \nl
      & \le \snorm{f} |x-y|+ |y| \snorm{f'} |x-y| \nl
      & \le \sqrt{\frac \pi2} |x-y|+2 |y| |x-y| \nl
      & = \left( \sqrt{\frac \pi2} + 2|y| \right) |x-y|
    \end{aligned}
  \end{comment} \nl
  &\le \E{ \left( \sqrt{\frac \pi2}+2 \abs{S_n^k} \right) \abs{S_n-S_n^k} } \nl
  &\le \E{ \left( \sqrt{\frac \pi2}+2 \abs{S_n^k} \right) \frac{|Y_k|}{\sqrt n} } \nl
  &\le \frac 1{\sqrt n} \left( \sqrt{\frac \pi2}+2 \E{\abs{S_n^k}} \right) \E{\abs{Y_k}} \nl
  &
  \begin{comment}
    \begin{aligned}
      \E{\abs{Y_k}}   & \le \sqrt{\E{Y_k^2}} = 1 \nl
      \E{\abs{S_n^k}} & \le \sqrt{\E{{S_n^k}^2}} \le \sqrt{\E{\frac{n}{n-1}{S_n^k}^2}} = 1
    \end{aligned}
  \end{comment} \nl
  &\le \frac 1{\sqrt n} \left( \sqrt{\frac \pi2}+2 \right)
\end{align}

\subsubsection{Estimation of the second term}

The second term can be estimated analogy to the first term:

\begin{align}
  \E{\abs{ \ft{ S_n^k + \Theta \frac{Y_k}{\sqrt n} } - \ft{ S_n^k } }} & \le \frac{ \E{\abs{Y_k}} }{\sqrt n} \left( \E{\Theta} \sqrt{\frac \pi2} + 2 \E{\Theta} \E{\abs{S_n^k}} \right) \nl
  &
  \begin{comment}
    \E{\Theta} = \frac 12 \text{ and } \E{\abs{S_n^k}} \le 1
  \end{comment} \nl
  & \le \frac{ \E{\abs{Y_k}} }{\sqrt n} \left(\frac 12\sqrt{\frac \pi2}+1\right) \nl
  &\begin{comment}
      \E{\abs{Y_k}} \le \sqrt{\E{Y_k^2}} = 1
  \end{comment} \nl
  & \le \frac{1}{\sqrt n} \left(\frac 12\sqrt{\frac \pi2}+1\right)
\end{align}

\subsubsection{Estimation of the third term}

\begin{align}
  & \E{\abs{ \h{S_n} - \h{ S_n^k }}} \nl
  &
  \begin{comment}
    \begin{aligned}
      h(y)-h(x) & = \int_y^x h'(s) \d{s} \nl
      &= (y-x) \int_0^1 h'(x+s(y-x)) \d{s} \nl
      &
      \begin{comment}
        \text{\parbox{0.45\textwidth}{
          Choose $\tilde \Theta$ uniformly distributed in $[0,1]$ and independent of all $Y_1,\ldots,Y_n$ and $\Theta$.
        }}
      \end{comment} \nl
      &= -\frac{y-x}{2\epsilon} \E{ \cf[z-\epsilon]{z+\epsilon}{x + \tilde \Theta (y-x)} }
    \end{aligned}
  \end{comment} \nl
  &= \frac{1}{2\epsilon} \E{ \abs{ S_n-S_n^k} \E{ \cf[z-\epsilon]{z+\epsilon}{S_n^k + \tilde \Theta \left( S_n-S_n^k \right)} }} \nl
  &= \frac{ \E{\abs{Y_k}} }{2\sqrt{n}\epsilon} \E{ \cf[z-\epsilon]{z+\epsilon}{S_n^k + \tilde \Theta \frac{Y_k}{\sqrt n} }} \nl
  &
  \begin{comment}
    \E{\abs{Y_k}} \le 1
  \end{comment} \nl
  &\le \frac{1}{2\sqrt{n}\epsilon} \E{ \cf[z-\epsilon]{z+\epsilon}{S_n^k + \tilde \Theta \frac{Y_k}{\sqrt n} }} \nl
  &= \frac{1}{2\sqrt{n}\epsilon} \P{z-\epsilon \le S_n^k + \tilde \Theta \frac{Y_k}{\sqrt n} \le z+\epsilon } \nl
  &
  \begin{comment}
    \tilde \Theta\frac{Y_k}{\sqrt n} \in \R
  \end{comment} \nl
  &\le \frac{1}{2\sqrt{n}\epsilon} \sup_{c\in\R} \P{z-c-\epsilon \le S_n^k \le z-c+\epsilon}
\end{align}

To estimate $\P{ a\le S_n^k\le b }$ we set $\tilde S_n^k = \tfrac{1}{\sqrt{n-1}} \sum_{j\neq k} Y_j$. $\tilde S_n^k$ is the standardized sum of all $Y_1,\ldots,Y_k$ without $Y_k$. We have $S_n^k = \sqrt{1-\tfrac 1n} \tilde S_n^k$. Thus

\begin{align}
  \P{a\le S_n^k\le b} &= \P{a \le \sqrt{1-\tfrac 1n} \tilde S_n^k \le b} \nl
  &= \P{\frac{a}{\sqrt{1-\tfrac 1n}} \le \tilde S_n^k \le \frac{b}{\sqrt{1-\tfrac 1n}}} \nl
  &
  \begin{multlined}
    = \P{\frac{a}{\sqrt{1-\tfrac 1n}} \le \tilde S_n^k \le \frac{b}{\sqrt{1-\tfrac 1n}}}- \P{\frac{a}{\sqrt{1-\tfrac 1n}} \le N \le \frac{b}{\sqrt{1-\tfrac 1n}}} \nl
    +\P{\frac{a}{\sqrt{1-\tfrac 1n}} \le N \le \frac{b}{\sqrt{1-\tfrac 1n}}}
  \end{multlined} \nl
  &
  \begin{comment}
    \P{\frac{a}{\sqrt{1-\tfrac 1n}} \le N \le \frac{b}{\sqrt{1-\tfrac 1n}}} \le \frac{b-a}{\sqrt{2\pi}\sqrt{1-\tfrac 1n}}
  \end{comment} \nl
  &
  \begin{multlined}
    \le \P{\frac{a}{\sqrt{1-\tfrac 1n}} \le \tilde S_n^k \le \frac{b}{\sqrt{1-\tfrac 1n}}}- \P{\frac{a}{\sqrt{1-\tfrac 1n}} \le N \le \frac{b}{\sqrt{1-\tfrac 1n}}} \nl
    +\frac{b-a}{\sqrt{2\pi}\sqrt{1-\tfrac 1n}}
  \end{multlined} \nl
  &
  \begin{comment}
    \begin{aligned}
      & \P{x \le \tilde S_n^k \le y} - \P{x\le N \le y} \nl
      &
      \begin{multlined}
        = \P{\tilde S_n^k \le y} - \P{\tilde S_n^k < x} \nl
        - \P{N \le y} + \P{N < x}
      \end{multlined} \nl
      &
      \begin{multlined}
        = \P{\tilde S_n^k \le y} - \P{N \le y} \nl
        + \P{N < x} - \P{\tilde S_n^k < x}
      \end{multlined} \nl
      &\le 2 \sup_{z\in\R} \abs{\P{\tilde S_n^k \le z}-\Phi(z)} \nl
      &= \frac{2 C_{n-1} \E{\abs{Y_1}^3}}{\sqrt{n-1}}
     \end{aligned}
  \end{comment} \nl
  & \le \frac{2 C_{n-1} \E{\abs{Y_1}^3}}{\sqrt{n-1}} + \frac{b-a}{\sqrt{2\pi}\sqrt{1-\tfrac 1n}}
\end{align}

\noindent Thus

\begin{align}
  \E{\abs{h(S_n) - h(N)}} & \le \frac{1}{2\sqrt{n}\epsilon} \sup_{c\in\R} \P{ z-c-\epsilon \le S_n^k \le z-c+\epsilon} \nl
  &\le \frac{2C_{n-1} \E{\abs{Y_1}^3}}{2\sqrt n\sqrt{n-1}\epsilon}+\frac{2\epsilon}{2\epsilon\sqrt{2\pi}\sqrt{n-1}} \nl
  &= \frac{C_{n-1} \E{\abs{Y_1}^3}}{\sqrt n\sqrt{n-1}\epsilon}+\frac{1}{\sqrt{2\pi}\sqrt{n-1}} \nl
\end{align}

\subsubsection{Estimation of the fourth term}

Also the fourth term can be estimated analogy to the third term:

\begin{align}
  \E{\abs{ \h{ S_n^k+\Theta\frac{Y_k^2}{n} } - \h{ S_n^k } }}
  &= \frac{1}{2\sqrt n \epsilon} \E{ \abs{Y_k} \abs{\Theta} \E{ \cf[z-\epsilon]{z+\epsilon}{S_n^k + \tilde \Theta\Theta\frac{Y_k}{\sqrt n}} } } \nl
  &
  \begin{comment}
    \E{\abs{Y_k}} \le 1 \text{ and } \E{\Theta} = \frac 12
  \end{comment} \nl
  &= \frac{1}{4\sqrt n \epsilon} \E{ \cf[z-\epsilon]{z+\epsilon}{S_n^k + \tilde \Theta\Theta\frac{Y_k}{\sqrt n}} } \nl
  & \le \frac{1}{4\sqrt n \epsilon} \sup_{c\in\R} \P{z-c-\epsilon \le S_n^k \le z-c+\epsilon} \nl
  &\le \frac{1}{2\sqrt{2\pi}\sqrt{n-1}} + \frac{C_{n-1} \E{\abs{Y_1}^3} }{2\sqrt  n\sqrt{n-1}\epsilon}
\end{align}

\subsection{Derivation of recursive relation for $C_n$}

\begin{align}
  \abs{\E{h(S_n)}-\E{h(N)}}
  &\le \frac 1n \sum_{k=1}^n \left(
  \begin{multlined}
      \E{\abs{\ft{S_n} - \ft{S_n^k}}} \nl
      +\E{\abs{\ft{ S_n^k +\Theta \frac{Y_k}{\sqrt n} } - \ft{ S_n^k }}} \nl
      +\E{\abs{\h{S_n} - \h{ S_n^k }}} \nl
      +\E{\abs{\h{ S_n^k +\Theta \frac{Y_k}{\sqrt n} } - \h{ S_n^k }}}
    \end{multlined}
  \right) \bnl
  &\le \frac 1n \sum_{k=1}^n \left(
    \begin{multlined}
      \left(\sqrt{\frac\pi2}+2\right) \frac1{\sqrt n} +\left(\frac 12\sqrt{\frac \pi2}+1\right)\frac{1}{\sqrt n}\nl
      +\frac{1}{\sqrt{2\pi}\sqrt{n-1}}+\frac{C_{n-1} \E{\abs{Y_1}^3} }{\sqrt{n}\sqrt{n-1}\epsilon}\nl
      +\frac{1}{2\sqrt{2\pi}\sqrt{n-1}}+\frac{C_{n-1} \E{\abs{Y_1}^3} }{2\sqrt{n}\sqrt{n-1}\epsilon}
    \end{multlined}
  \right) \bnl
  &
  \begin{multlined}
    = \left(\sqrt{\frac\pi2}+2\right) \frac1{\sqrt n} +\left(\frac 12\sqrt{\frac \pi2}+1\right)\frac{1}{\sqrt n}\nl
    +\frac{1}{\sqrt{2\pi}\sqrt{n-1}}+\frac{C_{n-1} \E{\abs{Y_1}^3} }{\sqrt{n}\sqrt{n-1}\epsilon}\nl
    +\frac{1}{2\sqrt{2\pi}\sqrt{n-1}}+\frac{C_{n-1} \E{\abs{Y_1}^3} }{2\sqrt{n}\sqrt{n-1}\epsilon}
  \end{multlined} \bnl
  & = \frac{1}{\sqrt{n}} \left(
  \begin{multlined}
    \sqrt{\frac\pi2} + 2 + \frac 12 \sqrt{\frac \pi2} + 1 +\frac{3\sqrt n}{2\sqrt{2\pi}\sqrt{n-1}} \nl
    + \frac{3C_{n-1} \E{\abs{Y_1}^3} }{2\sqrt{n-1}\epsilon}
  \end{multlined}
  \right) \bnl
  &
  \begin{comment}
    \E{\abs{Y_1}^3} \ge \E{Y_1^2}^{\tfrac 32} = 1
  \end{comment} \bnl
  & = \frac{ \E{\abs{Y_1}^3}  }{\sqrt{n}} \left(
  \begin{multlined}
    \sqrt{\frac\pi2} + 2 + \frac 12 \sqrt{\frac \pi2} + 1 +\frac{3\sqrt n}{2\sqrt{2\pi}\sqrt{n-1}} \nl
    + \frac{3C_{n-1} \E{\abs{Y_1}^3} }{2\sqrt{n-1}\epsilon}
  \end{multlined}
  \right) \bnl
  &
  \begin{comment}
    n \ge 2 \implies n-1 \ge \frac n2 \implies \sqrt{n-1} \ge \frac{\sqrt{n}}{\sqrt 2} 
  \end{comment} \bnl
  & = \frac{ \E{\abs{Y_1}^3}  }{\sqrt{n}} \left(
  \begin{multlined}
    \sqrt{\frac\pi2} + 2 + \frac 12 \sqrt{\frac \pi2} + 1 +\frac{3\sqrt 2}{2\sqrt{2\pi}} \nl
    + \frac{3 \sqrt{2} C_{n-1} \E{\abs{Y_1}^3} }{2\sqrt{n}\epsilon}
  \end{multlined}
  \right) \bnl
  &
  \begin{comment}
    \frac{\sqrt 2}2 \le 1
  \end{comment} \bnl
  & = \frac{ \E{\abs{Y_1}^3}  }{\sqrt{n}} \left( 5.726\ldots + \frac{3 C_{n-1} \E{\abs{Y_1}^3} }{\sqrt{n}\epsilon}
  \right) \bnl
  &\le \frac{6 \E{\abs{Y_1}^3} }{\sqrt n} + \frac{3 C_{n-1} \E{\abs{Y_1}^3}^2}{n\epsilon}
\end{align}

\noindent This inequality can be used in \eqref{berryproof:before_stein}:

\begin{align}
  \sup_{x\in\R} \abs{\P{S_n \le x} - \P{N \le x}} & \le \sup_{x\in\R} \Big\{ \abs{\E{h_{x,\epsilon}(S_n)} - \E{h_{x,\epsilon}(N)}} \Big\} + \frac{4\epsilon}{\sqrt{2\pi}} \bnl
  &\le \frac{6 \E{\abs{Y_1}^3} }{\sqrt n} + \frac{3 C_{n-1} \E{\abs{Y_1}^3}^2}{n\epsilon}+ \frac{4\epsilon}{\sqrt{2\pi}}
\end{align}

\noindent This inequality holds for all $\epsilon > 0$. We choose $\epsilon = \sqrt{\frac{C_{n-1}}{n}} \E{\abs{Y_1}^3}$. So we get

\begin{align}
  \sup_{x\in\R} \abs{\P{S_n \le x} - \P{N \le x}} & \le \frac{6 \E{\abs{Y_1}^3} }{\sqrt n} + \frac{3 \sqrt{C_{n-1}} \E{\abs{Y_1}^3}^2}{ \sqrt n \epsilon } + \frac{4\sqrt{C_{n-1}} \E{\abs{Y_1}^3} }{\sqrt{2\pi}\sqrt{n}} \nl
  &= \frac{ \E{\abs{Y_1}^3} }{\sqrt n} \left( 6+\left(3+\frac{4}{\sqrt{2\pi}}\right)\sqrt{C_{n-1}} \right)
\end{align}

\noindent We have defined $C_n$ such that it is the smallest number fulfilling 

\begin{align}
  \sup_{x\in\R} \abs{\P{S_n \le x} - \P{N \le x}} \le \frac{ \E{\abs{Y_1}^3} }{\sqrt n} C_n
\end{align}

\noindent Therefore we can conclude for $n \ge 2$

\begin{align}
  C_n \le 6 + \left(3 + \frac{4}{\sqrt{2\pi}}\right) \sqrt{C_{n-1}}
\end{align}

\subsection{Final induction}

No we proof by induction that $C_n < 33$. First note that

\begin{align}
  C_n & = \frac{\sqrt{n}}{\E{\abs{X_1}^3}} \sup_{x\in\R} \abs{\P{S_n \le x}-\Phi(x)} \nl
  &
  \begin{comment}
    \E{\abs{Y_1}^3} \ge 1
  \end{comment} \nl
  & \le \sqrt{n} \sup_{x\in\R} \abs{\P{S_n \le x}-\Phi(x)} \nl
  &
  \begin{comment}
    \begin{multlined}
      \P{S_n \le x} \in [0,1] \land \Phi(x) \in [0,1] \nl
      \implies \abs{\P{S_n \le x}-\Phi(x)} \le 1
    \end{multlined}
  \end{comment} \nl
  & \le \sqrt{n}
\end{align}

So we have $C_1 \le \sqrt{1} < 33$ and

\begin{align}
  C_{n+1} & \le 6+\left(3 + \frac{4}{\sqrt{2\pi}}\right) \sqrt{C_n} \nl
  & \le 6+\left(3 + \frac{4}{\sqrt{2\pi}}\right) \sqrt{33} \nl
  & = 32.40\ldots < 33
\end{align}

This proves $\sup_{n\in\N} C_n < \infty$ from which the Berry-Essen theorem follows.

    \chapter{Source code for plots}

\newcommand*{\src}[1]{\section{Figure \ref{#1}}

\lstinputlisting[language=Python]{"plots/#1.py"}}

\src{pmf}
\src{cdf}
\src{ml_pmf}
\src{ml_cdf}


    \printbibliography
    \markboth{}{}

    \backmatter
    
    \cleardoublepage
\phantomsection \label{acknowledgment}
\addcontentsline{toc}{chapter}{\protect Acknowledgment}

\chapter*{Acknowledgment}

I would like to thank my supervisor \href{http://www.mathematik.uni-muenchen.de/personen/professoren/pickl/index.html}{Prof. Dr. Peter Pickl} for his ideas and help throughout writing this thesis. I especially want to thank him for his great support and the freedom I had.


    \markboth{Acknowledgment}{Acknowledgment}
    \cleardoublepage

    \listoffigures
    \markboth{List of Figures}{List of Figures}
    \cleardoublepage

    \markboth{Statement of Authorship}{Statement of Authorship}
    \cleardoublepage
\phantomsection \label{statement_of_authorship}
\addcontentsline{toc}{chapter}{\protect Statement of Authorship}

\chapter*{Statement of Authorship}

Hiermit versichere ich, dass ich die vorliegende Bachelorarbeit mit dem Titel

\begin{center}
\textbf{The de Moivre-Laplace theorem}
\end{center}

\noindent selbständig verfasst und keine anderen als die angegebenen Quellen und Hilfsmittel benutzt habe.

\vspace*{2cm}

\noindent\rule{0.3\textwidth}{0.4pt}

\noindent Stephan Kulla, \newline 
München, den 15.02.2013

    \cleardoublepage
  \end{appendices}
\end{document}
